% =============================================================================
%  Fundamentos da linguagem
%  Cobre a Etapa 1 da trilha: syntax, casts, value_semantics, operators,
%  smart-pointers, oop, errors -- mais file_io, inserido depois de
%  smart-pointers porque um stream e o primeiro recurso RAII com estado.
% =============================================================================

% =============================================================================
\section{Fundamentos da linguagem}
\label{sec:fundamentos}
% =============================================================================

Este é o capítulo mais longo do relatório, e o mais rentável: quase tudo o que
vem depois --- containers, \emph{templates}, padrões de projeto, concorrência,
integração com bibliotecas --- é aplicação do que está aqui. O assunto comum às
oito seções é um só: \textbf{o que o compilador faz com o seu código}, e onde
ele faz algo diferente do que a leitura ingênua sugere.

A ordem interna reproduz a Etapa~1 da trilha do repositório, com duas
inserções. \projeto{operators} entra logo depois de \projeto{value\_semantics},
porque sobrecarregar um operador corretamente depende de saber quando um
argumento é copiado e quando é movido. E \projeto{file\_io} entra logo depois de
\projeto{smart-pointers}, porque um \emph{stream} é o primeiro recurso
\gls{raii} que se encontra na prática --- e o único da biblioteca padrão que
sinaliza erro por \emph{bit de estado} em vez de exceção, o que faz dele a
ponte natural entre a \cref{sec:smart-pointers} e a \cref{sec:errors}.

\begin{porque}[por que começar pelo que o compilador cala]
As tabelas de armadilhas que abrem quase todas as seções deste capítulo não são
folclore. Elas têm em comum a propriedade que torna o C++ difícil de aprender
sozinho: \textbf{o programa compila, roda e dá a resposta certa} --- até o dia
em que não dá. Um destrutor vazio não gera \emph{warning}; um
\texttt{std::vector<Base>} não gera \emph{warning}; um \classe{const\_cast}
seguido de escrita não gera \emph{warning}. Quem só lê a referência da
linguagem nunca encontra esses casos, porque a referência descreve o que é
válido, não o que é ilusório.
\end{porque}

% =============================================================================
\subsection{O mapa da sintaxe}
\label{sec:syntax}
% =============================================================================

O subprojeto \projeto{syntax} é o mapa: toda construção da linguagem ---
declaração, seleção, iteração, função, tipo, classe, \emph{namespace},
\emph{template}, exceção --- num só executável, comentada. É o subprojeto que se
lê primeiro e ao qual se volta como referência; onde um tópico tem um
subprojeto dedicado, ele mostra a \emph{forma} e aponta para lá.

\subsubsection{Declarações e inicialização}

Toda declaração em C++ tem a mesma forma --- especificadores, declarador,
inicializador --- e o inicializador é onde mora a primeira armadilha. Existem
seis maneiras de inicializar uma variável, e elas não são equivalentes:

\begin{lstlisting}[style=bt, language=C++,
    caption={As seis formas de inicializar, e o que cada uma garante.}]
int a = 42;      // copy-init
int b(42);       // direct-init
int c{42};       // list-init      ← a preferida
int d = {42};    // copy-list-init
int e{};         // value-init     → ZERO garantido
int f;           // default-init   → LIXO indeterminado (ler é UB)
\end{lstlisting}

A forma com chaves venceu porque \textbf{proíbe \emph{narrowing}}: a conversão
que perde informação deixa de ser silenciosa e passa a ser erro de compilação.

\begin{lstlisting}[style=bt, language=C++,
    caption={A diferença de uma chave.}]
int t = 3.14159;   // compila, trunca em silêncio
int s{3.14159};    // erro de compilação: narrowing conversion
\end{lstlisting}

\begin{armadilha}[com containers, \texttt{\{\}} muda o \emph{significado}]
A preferência por chaves tem uma exceção que dói. Para os containers da
biblioteca padrão existe um construtor que recebe
\texttt{std::initializer\_list}, e ele \emph{vence} os demais na resolução de
sobrecarga:
\texttt{std::vector<int> v(3,0)} constrói três zeros;
\texttt{std::vector<int> v\{3,0\}} constrói dois elementos, 3 e 0. As duas
linhas compilam e nenhuma gera \emph{warning}.
\end{armadilha}

\paragraph{\texttt{const} e ponteiros.} A regra prática é ler a declaração da
direita para a esquerda; o que muda é qual dos dois --- o alvo ou o ponteiro ---
está congelado (\cref{tab:const-ponteiro}).

\begin{table}[htbp]
\centering
\caption{As três combinações de \texttt{const} com ponteiro.}
\label{tab:const-ponteiro}
\footnotesize
\begin{tabular}{@{}lcc@{}}
\toprule
\textbf{Declaração} & \texttt{*p = x} & \texttt{p = \&y} \\
\midrule
\texttt{const int* p}       & não & sim \\
\texttt{int* const p}       & sim & não \\
\texttt{const int* const p} & não & não \\
\bottomrule
\end{tabular}
\end{table}

\paragraph{Duração e visibilidade.} A \emph{storage class} decide quanto tempo
o objeto vive; o \emph{linkage} decide quem enxerga o nome. As duas noções são
independentes e costumam ser confundidas.

\begin{table}[htbp]
\centering
\caption{\emph{Linkage} das construções que mais aparecem em cabeçalhos.}
\label{tab:linkage}
\footnotesize
\begin{tabular}{@{}ll@{}}
\toprule
\textbf{Construção} & \textbf{\emph{Linkage}} \\
\midrule
\texttt{int x;} global                                   & externa \\
\texttt{static int x;} / \texttt{namespace \{ int x; \}} & interna \\
\texttt{inline int x = 1;} (C++17)                       & externa, uma instância \\
\texttt{const int x = 1;} global                         & \textbf{interna} --- a surpresa comum \\
\bottomrule
\end{tabular}
\end{table}

A última linha é a que pega: uma constante global em um cabeçalho não viola a
ODR porque cada unidade de tradução ganha a sua própria cópia --- o que também
significa que o endereço dela é diferente em cada \texttt{.cpp}. O
\texttt{inline} do C++17 é o que permite ter constante em cabeçalho com
\emph{uma} instância só.

\subsubsection{Controle de fluxo}

São seis construções --- \texttt{if}, \texttt{switch}, \texttt{while},
\texttt{do-while}, \texttt{for}, \texttt{range-for} --- e três saltos. O que o
C++ moderno acrescentou não foi construção nova: foi \textbf{escopo menor}.

\begin{lstlisting}[style=bt, language=C++,
    caption={O \emph{init-statement}: a variável morre onde deixa de fazer sentido.}]
if (auto it = m.find(k); it != m.end()) { /* usa it */ }      // C++17
for (auto f = fn(); const auto& x : v) { /* usa f e x */ }     // C++20
\end{lstlisting}

\texttt{if constexpr} é diferente de tudo o mais nessa lista: ele não escolhe um
caminho em tempo de execução, ele \textbf{descarta um galho na compilação}. O
galho descartado nem precisa ser válido para aquele tipo --- é o que torna
possível escrever uma função genérica que chama \texttt{.size()} só quando o
tipo tem \texttt{.size()}.

\begin{regra}[\texttt{switch} sobre \texttt{enum class} não leva \texttt{default}]
Com um \texttt{enum} escopado e sem \texttt{default}, acrescentar um enumerador
novo faz o compilador apontar todos os \texttt{switch} que deixaram de ser
exaustivos. Escrever \texttt{default:} desliga esse alarme para sempre. A
ausência do \texttt{default} é uma \emph{feature}, não um descuido --- e é a
forma mais barata que a linguagem oferece de verificação de exaustividade.
\end{regra}

\subsubsection{Funções}

A tabela de passagem de parâmetros (\cref{tab:passagem}) resolve a maioria das
decisões de assinatura. Ela reaparece, refinada, na \cref{sec:value-semantics},
quando a discussão incluir \emph{move}.

\begin{table}[htbp]
\centering
\caption{Como passar um parâmetro, e o que cada forma comunica.}
\label{tab:passagem}
\footnotesize
\begin{tabular}{@{}lcccl@{}}
\toprule
\textbf{Forma} & \textbf{Copia?} & \textbf{Modifica?} & \textbf{Aceita nulo?} & \textbf{Quando usar} \\
\midrule
\texttt{T}        & sim & não     & não & tipos baratos, \emph{sink} \\
\texttt{T\&}      & não & sim     & não & saída, entrada-e-saída \\
\texttt{const T\&}& não & não     & não & \textbf{leitura --- o padrão} \\
\texttt{T*}       & não & sim     & sim & opcional, API C \\
\texttt{T\&\&}    & não & consome & não & \emph{move}, \emph{forwarding} \\
\bottomrule
\end{tabular}
\end{table}

A resolução de sobrecarga tenta, nesta ordem: correspondência exata, promoção,
conversão padrão, conversão definida pelo usuário, reticências. E o
\textbf{tipo de retorno não entra na assinatura} --- \texttt{int f(int)} e
\texttt{bool f(int)} colidem.

As variádicas têm duas formas, e a distância entre elas mede o que o C++ ganhou
sobre o C:

\begin{lstlisting}[style=bt, language=C++,
    caption={Variádica à moda C e à moda C++.}]
int soma(int n, ...);                     // sem tipos, contagem manual, UB fácil

template <typename... A>
auto soma(A... a) { return (a + ...); }   // fold expression: tipado, compile-time
\end{lstlisting}

\subsubsection{Tipos definidos pelo usuário}

\texttt{struct} e \texttt{class} são a mesma construção: muda só o acesso
padrão. As distinções que importam de verdade estão em outro lugar
(\cref{tab:enum,tab:union}).

\begin{table}[htbp]
\centering
\caption{\texttt{enum} contra \texttt{enum class}.}
\label{tab:enum}
\footnotesize
\begin{tabular}{@{}lll@{}}
\toprule
 & \texttt{enum} & \texttt{enum class} \\
\midrule
escopo               & vaza para fora        & exige \texttt{Cor::Verde} \\
converte para \texttt{int} & sozinho         & só com \classe{static\_cast} \\
tipo subjacente      & escolha do compilador & sua (\texttt{: uint8\_t}) \\
colisão de nomes     & acontece              & impossível \\
\bottomrule
\end{tabular}
\end{table}

\begin{table}[htbp]
\centering
\caption{\texttt{union} contra \texttt{std::variant}.}
\label{tab:union}
\footnotesize
\begin{tabular}{@{}lll@{}}
\toprule
 & \texttt{union} & \texttt{std::variant} \\
\midrule
quem guarda a etiqueta   & você                & ele \\
ler o membro errado      & \gls{ub} silencioso & \texttt{std::bad\_variant\_access} \\
tipos não triviais       & não                 & sim \\
\bottomrule
\end{tabular}
\end{table}

Do C++20 vieram os \emph{designated initializers}, a resposta tardia à falta de
argumentos nomeados: \texttt{Config c\{.host = "api", .porta = 443, .tls =
true\}}. A restrição a conhecer: a ordem dos campos escritos tem de seguir a
ordem de \emph{declaração} dos membros --- o C é mais permissivo aqui, e código
portado de C pode não compilar.

\subsubsection{Classes: a lista de inicialização não é açúcar sintático}

\begin{lstlisting}[style=bt, language=C++,
    caption={Inicializar e atribuir não são a mesma coisa.}]
Classe(const T& r) : membro(r)  { }        // 1 construção
Classe(const T& r) { membro = r; }         // 1 construção padrão + 1 atribuição
\end{lstlisting}

A segunda forma exige que o membro seja construível por padrão, e faz o dobro
do trabalho. Para membros \texttt{const} ou de tipo referência ela nem é uma
opção: eles só podem ser \emph{inicializados}.

\begin{armadilha}[a ordem da lista escrita é ignorada]
Os membros são construídos na ordem de \textbf{declaração} na classe, não na
ordem em que aparecem na lista de inicialização. Escrever a lista fora de ordem
compila, roda na ordem de declaração mesmo assim, e --- se um membro
inicializar a partir de outro --- lê um membro ainda não construído. O GCC
avisa com \texttt{-Wreorder}, que não vem em \texttt{-Wall}.
\end{armadilha}

\subsubsection{Organização: \emph{namespaces}, ADL e o pré-processador}

A busca de nome do C++ considera também os \emph{namespaces} dos
\emph{argumentos} da chamada --- é a \emph{Argument-Dependent Lookup}, e é o que
faz \texttt{std::cout << x} e o idioma \texttt{using std::swap; swap(a,b)}
funcionarem sem qualificação explícita. É também o mecanismo que torna os
\emph{hidden friends} da \cref{sec:operators} eficientes.

Sobre macros, a regra é curta: quase toda uma tem substituto melhor
(\cref{tab:macros}). Sobram os casos que \emph{exigem} manipulação de texto ---
guardas de inclusão, \texttt{\_\_LINE\_\_}, portabilidade condicional.

\begin{table}[htbp]
\centering
\caption{O que usar no lugar de cada macro.}
\label{tab:macros}
\footnotesize
\begin{tabular}{@{}ll@{}}
\toprule
\textbf{Macro} & \textbf{Substituto} \\
\midrule
\texttt{\#define PI 3.14}   & \texttt{constexpr double pi = 3.14;} \\
\texttt{\#define MAX(a,b)}  & \texttt{template <class T> T max(T, T)} \\
\texttt{\#define DEBUG\_ON} & \texttt{constexpr bool} + \texttt{if constexpr} \\
\bottomrule
\end{tabular}
\end{table}

\begin{regra}[nunca \texttt{using namespace} no escopo global de um cabeçalho]
A diretiva impõe a decisão a \emph{todo} arquivo que incluir aquele cabeçalho,
direta ou indiretamente, e a consequência típica --- uma ambiguidade nova que
aparece em um \texttt{.cpp} que ninguém tocou --- é praticamente impossível de
diagnosticar a partir da mensagem de erro.
\end{regra}

\subsubsection{Expressões: precedência não é ordem de avaliação}

Duas armadilhas de expressão aparecem em produção com frequência
desproporcional ao quanto são discutidas.

A primeira é a comparação com sinais misturados. Numa comparação entre
\texttt{int} e \texttt{unsigned}, o \texttt{int} é convertido --- e $-1$ vira
$4\,294\,967\,295$:

\begin{lstlisting}[style=bt, language=C++,
    caption={A comparação que inverte de resultado.}]
int a = -1;  unsigned b = 1;
a < b        // FALSE
\end{lstlisting}

É daí que vem o \texttt{-Wsign-compare} ao escrever \texttt{int i; i <
v.size()}, e é por isso que o índice de um laço sobre container deve ser
\texttt{std::size\_t} --- ou, melhor, não existir, substituído por um
\emph{range-for}.

A segunda é confundir precedência com ordem de avaliação. \texttt{1 << 2 + 3} é
\texttt{1 << 5}, porque \texttt{+} tem precedência maior que \texttt{<<}. E,
independentemente disso, a ordem em que os operandos de uma expressão são
avaliados é, em geral, não especificada.

\subsubsection{As armadilhas silenciosas, reunidas}

A \cref{tab:armadilhas-syntax} recolhe o núcleo do argumento do subprojeto
\projeto{syntax}: cada linha é uma construção que compila, e a última coluna diz
se você fica sabendo.

\begin{table}[htbp]
\centering
\caption{Construções que não significam o que parecem --- e se o compilador
         avisa. Cada uma tem demonstração executável em \texttt{src/syntax/}.}
\label{tab:armadilhas-syntax}
\footnotesize
\begin{tabular}{@{}llc@{}}
\toprule
\textbf{Escrito assim} & \textbf{O que acontece} & \textbf{Avisa?} \\
\midrule
\texttt{Widget w();}                  & declara uma \emph{função}          & não \\
\texttt{int x;} local                 & lixo indeterminado; ler é \gls{ub} & não \\
\texttt{if (a) f(); g();}             & \texttt{g()} roda sempre           & às vezes \\
\texttt{052}                          & é octal, vale 42                   & não \\
\texttt{\#define Q(x) x*x}, \texttt{Q(1+2)} & vira \texttt{1+2*1+2} = 5    & não \\
\texttt{a < b}, \texttt{int} e \texttt{unsigned} & $-1$ vira $4\,294\,967\,295$ & \texttt{-Wsign-compare} \\
\texttt{int t = 3.99;}                & trunca em silêncio                 & não \\
\texttt{delete p} para \texttt{new T[n]} & \gls{ub}                        & não \\
\texttt{[=]} num método               & captura \texttt{this} por ponteiro & não \\
\texttt{case} sem \texttt{break}      & cai para o seguinte                & \texttt{-Wimplicit-fallthrough} \\
\bottomrule
\end{tabular}
\end{table}

% =============================================================================
\subsection{O sistema de tipos e as conversões explícitas}
\label{sec:casts}
% =============================================================================

O subprojeto \projeto{casts} é o menor pré-requisito da trilha, e ensina uma
coisa que a sintaxe esconde: os quatro \emph{casts} do C++ \textbf{não são
quatro sinônimos com nomes diferentes}. Cada um é uma operação distinta, com
verificações distintas e riscos distintos --- e o \emph{cast} à moda C é o
apagador dessa distinção.

\subsubsection{Os quatro, em ordem de segurança}

\begin{table}[htbp]
\centering
\caption{Os quatro \emph{casts} do C++: o que fazem e o que verificam.}
\label{tab:casts}
\footnotesize
\begin{tabular}{@{}lll@{}}
\toprule
\textbf{Cast} & \textbf{Faz} & \textbf{Verifica} \\
\midrule
\texttt{static\_cast}      & conversões relacionadas   & compilação \\
\classe{dynamic\_cast}     & desce na hierarquia       & \textbf{execução} (exige \texttt{virtual}) \\
\texttt{const\_cast}       & remove/adiciona \texttt{const} & nada --- escrever pode ser \gls{ub} \\
\classe{reinterpret\_cast} & reinterpreta os bits      & nada \\
\bottomrule
\end{tabular}
\end{table}

\paragraph{\texttt{static\_cast} --- o padrão.} Pede ao compilador que confira,
\emph{em compilação}, se a conversão faz sentido. Se não fizer, o programa não
compila --- que é o comportamento desejável. Serve para conversões numéricas
explícitas, para subir na hierarquia (sempre seguro) e para descer quando você
\emph{sabe} o tipo real. Não serve entre tipos não relacionados, e é justamente
por recusá-los que ele é o \emph{cast} padrão.

\paragraph{\texttt{dynamic\_cast} --- descer com verificação.} Usa \gls{rtti}
para conferir em tempo de execução se o objeto realmente é do tipo pedido. Exige
que a base tenha ao menos uma função virtual, e falha de duas maneiras
diferentes conforme a forma usada: com ponteiro devolve \texttt{nullptr}; com
referência lança \classe{std::bad\_cast}.

\begin{lstlisting}[style=bt, language=C++,
    caption={As duas formas de falha do \texttt{dynamic\_cast}.}]
Animal* animal = new Cachorro();

if (Cachorro* c = dynamic_cast<Cachorro*>(animal)) {   // nullptr se falhar
    c->buscar();
}

try {
    Cachorro& cc = dynamic_cast<Cachorro&>(*animal);   // lança se falhar
    cc.buscar();
} catch (const std::bad_cast& e) {
    fmt::print("não é um Cachorro: {}\n", e.what());
}
\end{lstlisting}

O custo não é desprezível: \texttt{dynamic\_cast} percorre a árvore de herança
em tempo de execução e não é $O(1)$. A \cref{sec:oop} traz a medição, e a
\cref{sec:oop-rtti} traz as alternativas --- que, quase sempre, se resumem a
uma função virtual que não foi escrita.

\paragraph{\texttt{const\_cast} --- a única aplicação legítima.} Remover
\texttt{const} é aceitável em exatamente um cenário: atravessar uma interface
antiga que esqueceu de declarar o parâmetro como \texttt{const}, sabendo que ela
não modifica o objeto.

\begin{armadilha}[modificar um objeto genuinamente \texttt{const} é sempre \gls{ub}]
\begin{lstlisting}[style=bt, language=C++, frame=none,
                   xleftmargin=0pt, xrightmargin=0pt]
const int x = 42;              // pode estar em memória somente-leitura
int* p = const_cast<int*>(&x);
*p = 100;                      // UB: pode falhar, pode ser ignorado,
                               // pode "funcionar" e mudar em -O2
\end{lstlisting}
A regra é sobre a \emph{origem} do objeto, não sobre o tipo do ponteiro: se ele
foi declarado \texttt{const}, escrever nele é \gls{ub} independentemente de
quantos \emph{casts} houver no caminho.
\end{armadilha}

\paragraph{\texttt{reinterpret\_cast} --- e por que \classe{std::bit\_cast} o
substitui.} Tratar os bytes de um tipo como se fossem outro é legítimo em
interoperação com C, em aritmética de endereços e ao inspecionar a
representação. O que não é legítimo é \emph{ler pelo tipo errado}: acessar um
\texttt{float} através de um \texttt{int*} viola \emph{strict aliasing}, e o
compilador tem permissão para supor que os dois ponteiros nunca apontam para a
mesma memória.

\begin{lstlisting}[style=bt, language=C++,
    caption={A forma errada e as duas formas corretas de ver os bits de um \texttt{float}.}]
float f = 1.0f;

int* i = reinterpret_cast<int*>(&f);
*i = 42;                                  // UB: viola strict aliasing

int bits;
std::memcpy(&bits, &f, sizeof(bits));     // correto, e o compilador otimiza
auto bits2 = std::bit_cast<std::uint32_t>(f);   // C++20: correto e constexpr
\end{lstlisting}

\subsubsection{Por que o \emph{cast} à moda C é o pior dos cinco}

\texttt{(Tipo)valor} tenta, em ordem, cinco combinações --- \texttt{static\_cast},
\texttt{const\_cast}, \texttt{reinterpret\_cast} e duas composições. Quem
escolhe é o compilador, e o programador não fica sabendo qual foi.

\begin{lstlisting}[style=bt, language=C++,
    caption={O mesmo \gls{ub}, uma vez escondido e outra vez visível.}]
const int x = 42;
int* p1 = (int*)&x;                 // faz const_cast em silêncio
int* p2 = const_cast<int*>(&x);     // a intenção fica escrita
\end{lstlisting}

Há um segundo argumento, mais prático que teórico: \texttt{static\_cast} é
\emph{pesquisável}. Auditar um projeto atrás de todas as conversões perigosas é
um \texttt{grep} quando os \emph{casts} têm nome, e é impossível quando eles são
um par de parênteses. O compilador ajuda a manter a disciplina com
\texttt{-Wold-style-cast}.

\begin{figure}[htbp]
\centering
\begin{tikzpicture}[node distance=5mm and 12mm]
  \node[caixaforte] (q) {preciso converter};
  \node[caixa, below left=8mm and 14mm of q]  (num) {conversão\\numérica};
  \node[caixa, below=of q]                    (hier) {desce na\\hierarquia};
  \node[caixa, below right=8mm and 14mm of q] (bits) {reinterpretar\\bytes};
  \node[caixaok, below=of num]  (sc) {\texttt{static\_cast}};
  \node[caixa, below=of hier]   (sabe) {sei o tipo\\real?};
  \node[caixaok, below=of bits] (bc) {\texttt{bit\_cast}\\ou \texttt{memcpy}};
  \node[caixaok, below left=6mm and 2mm of sabe]  (sc2) {\texttt{static\_cast}};
  \node[caixa, below right=6mm and 2mm of sabe]   (dc) {\texttt{dynamic\_cast}};
  \draw[seta] (q) -- (num);  \draw[seta] (q) -- (hier);  \draw[seta] (q) -- (bits);
  \draw[seta] (num) -- (sc); \draw[seta] (hier) -- (sabe); \draw[seta] (bits) -- (bc);
  \draw[seta] (sabe) -- node[rotulo,left] {sim} (sc2);
  \draw[seta] (sabe) -- node[rotulo,right] {não} (dc);
\end{tikzpicture}
\caption{A árvore de decisão de \texttt{src/casts/}. Nenhum galho termina em
         \emph{cast} à moda C, e apenas um termina em \texttt{dynamic\_cast} ---
         que, ao aparecer numa cadeia de \texttt{if}, quase sempre indica uma
         função virtual faltando.}
\label{fig:arvore-casts}
\end{figure}

% =============================================================================
\subsection{Semântica de valor: cópia, \emph{move} e elisão}
\label{sec:value-semantics}
% =============================================================================

O subprojeto \projeto{value\_semantics} sustenta todos os outros. \classe{std::move}
aparece em dezenas de arquivos do repositório; é aqui que se explica o que ele
realmente faz --- um \texttt{static\_cast}, e nada mais.

O que torna esta seção urgente é a natureza dos erros que ela previne: nenhum
deles gera \emph{warning}, e todos custam desempenho ou correção
(\cref{tab:erros-valor}). O programa do subprojeto roda com contadores
instrumentados, de modo que cada afirmação da tabela é \emph{medida}, não
alegada.

\begin{table}[htbp]
\centering
\caption{Cinco erros de semântica de valor. Nenhum é apontado pelo compilador
         com os avisos padrão.}
\label{tab:erros-valor}
\footnotesize
\begin{tabular}{@{}ll@{}}
\toprule
\textbf{Erro} & \textbf{Consequência} \\
\midrule
declarar \texttt{\textasciitilde T() \{\}} vazio        & suprime os \emph{moves} --- todo ``\emph{move}'' vira cópia \\
esquecer \texttt{noexcept} no \emph{move ctor}          & \texttt{std::vector} passa a copiar ao realocar \\
\texttt{return std::move(x)}                            & desliga o NRVO: de zero operações para uma \\
\texttt{std::move} onde cabia \classe{std::forward}     & rouba o objeto do chamador sem avisar \\
repassar \texttt{T\&\&} sem \texttt{forward}            & vira \emph{lvalue} no caminho; o \emph{move} se perde \\
\bottomrule
\end{tabular}
\end{table}

\subsubsection{Categorias de valor}

Toda expressão tem um tipo \textbf{e} uma \gls{categoriadevalor}. A categoria
responde a duas perguntas independentes: \emph{tem identidade?} e \emph{pode ser
movida?} (\cref{fig:categorias}).

\begin{figure}[htbp]
\centering
\begin{tikzpicture}[node distance=6mm and 6mm]
  \node[caixaforte] (e) {expressão};
  \node[caixa, below left=8mm and 6mm of e]  (gl) {glvalue\\\scriptsize tem identidade};
  \node[caixa, below right=8mm and 6mm of e] (rv) {rvalue\\\scriptsize pode ser movida};
  \node[folha, below=8mm of gl, xshift=-14mm] (lv) {\textbf{lvalue}\\\scriptsize \texttt{x}, \texttt{*p}, \texttt{arr[0]}};
  \node[folha, below=8mm of e]                (xv) {\textbf{xvalue}\\\scriptsize \texttt{std::move(x)}};
  \node[folha, below=8mm of rv, xshift=14mm]  (pv) {\textbf{prvalue}\\\scriptsize \texttt{42}, \texttt{f()}, \texttt{T\{\}}};
  \draw[aresta] (e) -- (gl); \draw[aresta] (e) -- (rv);
  \draw[aresta] (gl) -- (lv); \draw[aresta] (gl) -- (xv);
  \draw[aresta] (rv) -- (xv); \draw[aresta] (rv) -- (pv);
\end{tikzpicture}
\caption{As categorias de valor. O \emph{xvalue} é o único que pertence aos dois
         ramos: tem identidade \emph{e} pode ser movido --- é exatamente o que
         \texttt{std::move} produz.}
\label{fig:categorias}
\end{figure}

\begin{armadilha}[tudo o que tem nome é \emph{lvalue}]
\begin{lstlisting}[style=bt, language=C++, frame=none,
                   xleftmargin=0pt, xrightmargin=0pt]
void recebe(std::string&& s) {
    alvo(s);              // → escolhe o overload de LVALUE
    alvo(std::move(s));   // → escolhe o overload de rvalue
}
\end{lstlisting}
O parâmetro \texttt{s} foi \emph{declarado} como referência a \emph{rvalue},
mas a \emph{expressão} \texttt{s} tem nome --- e tudo o que tem nome é
\emph{lvalue}. É por isso que repassar um \texttt{T\&\&} sem
\texttt{std::move}/\texttt{std::forward} produz uma cópia. A regra existe por
segurança: sem ela, usar o parâmetro duas vezes esvaziaria o objeto na primeira.
\end{armadilha}

\paragraph{\texttt{std::move} é um \emph{cast}.} A implementação cabe em três
linhas, e nenhum byte é copiado ao chamá-la:

\begin{lstlisting}[style=bt, language=C++,
    caption={O que \texttt{std::move} faz, escrito por extenso.}]
template <typename T>
constexpr std::remove_reference_t<T>&& meu_move(T&& t) noexcept {
    return static_cast<std::remove_reference_t<T>&&>(t);
}
\end{lstlisting}

O \gls{move} de fato só acontece quando alguém \emph{constrói ou atribui} a
partir do \emph{rvalue} resultante. \texttt{std::move} apenas autoriza.

\paragraph{\texttt{decltype} preserva a categoria.} O detalhe dos parênteses
surpreende, mas é consistente: \texttt{decltype(x)} pergunta pelo tipo
\emph{declarado} de \texttt{x}, enquanto \texttt{decltype((x))} pergunta pelo
tipo da \emph{expressão} \texttt{(x)}, que é um \emph{lvalue} e portanto
\texttt{int\&}.

\subsubsection{Rule of 0 / 3 / 5}

Um tipo tem seis operações que o compilador pode gerar sozinho: construtor
padrão, destrutor, construtor e atribuição de cópia, construtor e atribuição de
\emph{move}. As regras de geração são interligadas, e a consequência prática é
a armadilha mais cara do capítulo.

\begin{armadilha}[um destrutor vazio custa todos os \emph{moves}]
\begin{lstlisting}[style=bt, language=C++, frame=none,
                   xleftmargin=0pt, xrightmargin=0pt]
struct ComDestrutor {
    std::vector<int> dados;
    ~ComDestrutor() {}      // isto sozinho suprime os DOIS moves
};
\end{lstlisting}
Declarar o destrutor --- ainda que vazio --- impede a geração implícita do
construtor e da atribuição de \emph{move}. O código continua compilando e
rodando; só fica mais lento, porque todo ``\emph{move}'' cai silenciosamente na
cópia. Saída real do subprojeto:
\begin{lstlisting}[style=bt, language=bash, frame=none,
                   xleftmargin=0pt, xrightmargin=0pt]
ComDestrutor  — move real disponível? NÃO — cai silenciosamente na cópia
SemDestrutor  — move real disponível? sim
\end{lstlisting}
Note que \texttt{std::is\_move\_constructible\_v} responderia \texttt{true} nos
dois casos, porque \texttt{const T\&} também aceita \emph{rvalue}. O
discriminador honesto é \texttt{std::is\_nothrow\_move\_constructible\_v}.
\end{armadilha}

\paragraph{\texttt{noexcept} decide se o \texttt{vector} move ou copia.} Ao
realocar, \texttt{std::vector} precisa oferecer a garantia forte: se algo falhar
no meio da transferência, o vetor original tem de continuar íntegro. Um
construtor de \emph{move} que pode lançar não permite essa promessa --- depois
de mover metade dos elementos, não há como voltar atrás. Por isso o container
consulta \texttt{std::move\_if\_noexcept} e, na dúvida, copia:

\begin{lstlisting}[style=bt, language=bash,
    caption={Saída de \texttt{src/value\_semantics/}: a única diferença entre os
             dois tipos é a palavra \texttt{noexcept}.}]
MoveNoexcept    → 7 moves, 0 cópias
MovePodeLancar  → 0 moves, 7 cópias
\end{lstlisting}

\paragraph{\emph{Copy-and-swap}.} Uma única atribuição, recebida \textbf{por
valor}, cobre cópia e \emph{move} de uma vez --- e ganha de graça a garantia
forte e o tratamento correto da autoatribuição:

\begin{lstlisting}[style=bt, language=C++,
    caption={\emph{Copy-and-swap} em \texttt{src/value\_semantics/demo\_rule\_of\_five.hpp}
             (condensado).}]
class Cofre {
public:
    explicit Cofre(std::size_t n) : tamanho_(n), dados_(new int[n]{}) {}
    ~Cofre() { delete[] dados_; }

    Cofre(const Cofre& o) : tamanho_(o.tamanho_), dados_(new int[o.tamanho_]) {
        for (std::size_t i = 0; i < tamanho_; ++i) dados_[i] = o.dados_[i];
    }

    Cofre(Cofre&& o) noexcept : tamanho_(o.tamanho_), dados_(o.dados_) {
        o.dados_ = nullptr;
        o.tamanho_ = 0;
    }

    // A única atribuição declarada — cobre os dois casos.
    Cofre& operator=(Cofre outro) noexcept {
        trocar(*this, outro);
        return *this;
    }

    friend void trocar(Cofre& a, Cofre& b) noexcept {
        std::swap(a.tamanho_, b.tamanho_);
        std::swap(a.dados_,   b.dados_);
    }

private:
    std::size_t tamanho_ = 0;
    int*        dados_   = nullptr;
};
\end{lstlisting}

Chamada com \emph{lvalue}, o parâmetro é construído por cópia; com
\emph{rvalue}, por \emph{move}. Um corpo para os dois casos.

\begin{regra}[declare apenas a versão por valor]
Ter \texttt{operator=(T)} e \texttt{operator=(T\&\&)} ao mesmo tempo é ambíguo
para argumentos \emph{rvalue}, e o programa não compila. O \emph{copy-and-swap}
é uma escolha exclusiva, não uma adição.
\end{regra}

\subsubsection{Elisão de cópia}

Desde o C++17, retornar um \emph{prvalue} não cria temporário nenhum: o objeto é
construído diretamente no destino. Isso deixou de ser otimização e passou a ser
\emph{regra da linguagem} --- e a prova é um tipo que não pode ser copiado nem
movido:

\begin{lstlisting}[style=bt, language=C++,
    caption={A elisão obrigatória, em
             \texttt{src/value\_semantics/demo\_copy\_elision.hpp}.}]
struct Imovel {
    int v;
    explicit Imovel(int x) : v(x) {}
    Imovel(const Imovel&)            = delete;
    Imovel(Imovel&&)                 = delete;
    Imovel& operator=(const Imovel&) = delete;
    Imovel& operator=(Imovel&&)      = delete;
};
Imovel faz_imovel() {
    return Imovel(42);          // compila: nenhuma cópia/move é sequer considerada
}
\end{lstlisting}

\begin{armadilha}[\texttt{return std::move(local)} é uma pessimização]
A NRVO --- elisão para uma variável local nomeada --- continua sendo opcional e
depende de o compilador enxergar o \emph{nome} da variável no \texttt{return}.
Envolvê-lo em \texttt{std::move} produz um \emph{xvalue}, que já não é um nome:
a elisão deixa de se aplicar e o que era zero operação vira um \emph{move}.
\begin{center}\footnotesize
\begin{tabular}{@{}lcc@{}}
\toprule
\textbf{Versão} & \textbf{Cópias} & \textbf{\emph{Moves}} \\
\midrule
\texttt{return local;} (NRVO)     & 0 & 0 \\
\texttt{return std::move(local);} & 0 & \textbf{1} \\
\bottomrule
\end{tabular}
\end{center}
O GCC avisa com \texttt{-Wpessimizing-move}. Retornar uma variável local já
trata a expressão como \emph{rvalue} quando a elisão não se aplica --- por
exemplo quando há dois candidatos em ramos diferentes --- de modo que o
\emph{fallback} natural já é o \emph{move}, sem ajuda.
\end{armadilha}

\subsubsection{\emph{Perfect forwarding}}

Repassar um argumento preservando tipo, constância e categoria exige duas peças
que só funcionam \emph{juntas}: o parâmetro \texttt{T\&\&} deduzido e o
\texttt{std::forward<T>} no repasse.

\begin{lstlisting}[style=bt, language=C++,
    caption={As duas peças do \gls{forwarding}.}]
template <typename T>
void repassa(T&& arg) {          // 1. referência universal (T é DEDUZIDO)
    alvo(std::forward<T>(arg));  // 2. forward condicional
}
\end{lstlisting}

O mecanismo é o \emph{colapso de referências}: quando \texttt{T} é deduzido em
\texttt{T\&\&}, um argumento \emph{lvalue} de tipo \texttt{X} deduz \texttt{T =
X\&} e o parâmetro colapsa para \texttt{X\&}; um \emph{rvalue} deduz \texttt{T =
X} e o parâmetro fica \texttt{X\&\&}. Das quatro combinações possíveis, apenas
\texttt{\&\& \&\&} sobrevive como \emph{rvalue}.

\begin{armadilha}[nem todo \texttt{T\&\&} é referência universal]
\begin{lstlisting}[style=bt, language=C++, frame=none,
                   xleftmargin=0pt, xrightmargin=0pt]
void f(std::string&&);                          // rvalue ref pura
template<class T> void f(std::vector<T>&&);     // rvalue ref pura: T não está
                                                // na posição do parâmetro
template<class T> void f(T&&);                  // referência universal
\end{lstlisting}
A forma só é universal quando o parâmetro é exatamente \texttt{T\&\&} com
\texttt{T} deduzido naquela posição. Nos outros dois casos, aplicar
\texttt{std::forward} não faz sentido --- e aplicar \texttt{std::move} está
correto.
\end{armadilha}

\begin{regra}[\texttt{move} em tipo concreto, \texttt{forward} em tipo deduzido]
\texttt{std::move} é incondicional: rouba sempre. \texttt{std::forward<T>} é
condicional: rouba só se o chamador ofereceu um \emph{rvalue}. Trocar um pelo
outro num \emph{template} esvazia silenciosamente o objeto de quem chamou ---
o programa continua correto para o primeiro uso e falha no segundo.
\end{regra}

\subsubsection{\emph{Move-only}, \emph{ref-qualifiers} e \emph{sink parameters}}

Alguns recursos não podem ser duplicados: descritor de arquivo, \emph{lock},
\emph{socket}, \emph{thread}. Para eles, copiar não é caro --- é \emph{errado},
e a forma de dizer isso é \texttt{= delete}, que torna o erro incompilável em
vez de detectável em produção.

Os \emph{ref-qualifiers} permitem que um método pergunte se \texttt{*this} é
descartável, e portanto se pode roubar o próprio conteúdo:

\begin{lstlisting}[style=bt, language=C++,
    caption={O mesmo método, duas implementações conforme a categoria de \texttt{*this}.}]
const std::string& texto() const&  { return texto_; }             // lvalue: empresta
std::string        texto() &&      { return std::move(texto_); }  // rvalue: rouba
\end{lstlisting}

Servem também para \emph{proibir}: declarar \texttt{T\& operator[](int) \&\& =
delete;} impede devolver uma referência para dentro de um temporário que morre
na mesma expressão --- um \emph{dangling} que, sem isso, compila.

Quando a função vai \textbf{guardar} uma cópia, a assinatura correta é por valor
seguida de \texttt{std::move} --- o \emph{sink parameter}:

\begin{lstlisting}[style=bt, language=C++,
    caption={Um \emph{overload} que atende \emph{lvalue} e \emph{rvalue} de
             forma ótima.}]
void adicionar(std::string nome) {          // por valor
    nomes_.push_back(std::move(nome));      // o move é obrigatório aqui
}
\end{lstlisting}

Com \emph{lvalue} há uma cópia --- inevitável, porque o chamador quer manter o
seu objeto; com \emph{rvalue} há um \emph{move}, que é barato. A alternativa de
escrever \texttt{const T\&} \emph{e} \texttt{T\&\&} explodiria em $2^n$
\emph{overloads} para $n$ parâmetros.

\begin{table}[htbp]
\centering
\caption{A tabela de assinaturas, agora completa.}
\label{tab:assinaturas}
\footnotesize
\begin{tabular}{@{}ll@{}}
\toprule
\textbf{Intenção} & \textbf{Assinatura} \\
\midrule
só ler o argumento                & \texttt{const T\&} \\
ler e modificar o do chamador     & \texttt{T\&} \\
guardar uma cópia                 & \texttt{T} por valor + \texttt{std::move} dentro \\
repassar em \emph{template}       & \texttt{T\&\&} + \texttt{std::forward<T>} \\
tipo pequeno e trivial            & \texttt{T} por valor \\
\bottomrule
\end{tabular}
\end{table}

% =============================================================================
\subsection{Sobrecarga de operadores}
\label{sec:operators}
% =============================================================================

O subprojeto \projeto{operators} responde à pergunta de quando fazer um tipo se
comportar como um tipo embutido --- e, mais importante, quando não fazer. A
regra que resolve a maioria das decisões é uma só:

\begin{regra}[não surpreenda]
Se \texttt{+} não for comutativo, associativo e barato, escreva um método com
nome. \texttt{a.concatenar(b)} nunca confundiu ninguém; \texttt{a + b} com
semântica inesperada confunde todo mundo, inclusive o autor, seis meses depois.
\end{regra}

\subsubsection{Membro ou função livre}

A diferença prática está na \textbf{conversão implícita do operando da
esquerda}. Um operador membro não converte o próprio \texttt{this}:

\begin{lstlisting}[style=bt, language=C++,
    caption={A assimetria que decide a forma.}]
m + 5.0   // funciona com operador membro (o double vira Dinheiro)
5.0 + m   // NÃO COMPILA: o `this` não converte
\end{lstlisting}

Daí a divisão da \cref{tab:forma-operador}. O caso dos operadores de
\emph{stream} não é escolha: \texttt{operator<<} tem de ser livre porque o
operando da esquerda é o \texttt{std::ostream}, que não é seu.

\begin{table}[htbp]
\centering
\caption{Qual forma usar para cada família de operadores.}
\label{tab:forma-operador}
\footnotesize
\begin{tabular}{@{}ll@{}}
\toprule
\textbf{Operador} & \textbf{Forma} \\
\midrule
simétrico (\texttt{+ - * / ==})      & função \textbf{livre} \\
modifica o objeto (\texttt{+=})      & \textbf{membro} \\
\texttt{= [] () -> ->*}              & \textbf{membro} (exigido pela linguagem) \\
\texttt{<< >>} de \emph{stream}      & \textbf{livre} (exigido) \\
\bottomrule
\end{tabular}
\end{table}

O idioma que reconcilia as duas primeiras linhas é a \textbf{forma canônica}:
a lógica vive no membro que modifica, e o operador simétrico é um
\emph{hidden friend} que a reaproveita.

\begin{lstlisting}[style=bt, language=C++,
    caption={A forma canônica --- e o motivo de \texttt{a} vir por valor.}]
Tipo& operator+=(const Tipo& o);                    // membro: aqui mora a lógica
friend Tipo operator+(Tipo a, const Tipo& b)        // livre: `a` por VALOR
    { return a += b; }
\end{lstlisting}

O parâmetro \texttt{a} por valor é a cópia de trabalho --- e, se o chamador
passou um temporário, ele é \emph{movido} para lá em vez de copiado
(\cref{sec:value-semantics}). Definir o \texttt{friend} \emph{dentro} do corpo
da classe tem três ganhos: o \emph{namespace} não é poluído, a resolução de
sobrecarga considera menos candidatos --- o que se nota no tempo de compilação
de projetos grandes --- e os dois operandos continuam podendo converter.

\begin{armadilha}[quatro operadores que não devem ser sobrecarregados]
\texttt{\&\&} e \texttt{||} sobrecarregados \textbf{perdem o curto-circuito}:
ambos os lados passam a ser avaliados, sempre. \texttt{,} perde a ordem de
avaliação garantida. \texttt{\&} unário quebra código genérico que espera obter
o endereço. Os três casos compilam sem reclamação e mudam o significado de
expressões escritas por terceiros.
\end{armadilha}

\subsubsection{\classe{operator<=>}: seis operadores viram um}

\begin{lstlisting}[style=bt, language=C++,
    caption={Comparação total em uma linha, em
             \texttt{src/operators/demo\_comparacao.hpp}.}]
struct Versao {
    int maior, menor, patch;

    // UMA linha. O compilador compara os membros na ordem de DECLARAÇÃO e
    // gera <, <=, >, >=, == e != a partir daqui.
    auto operator<=>(const Versao&) const = default;
};
\end{lstlisting}

A comparação gerada é lexicográfica na ordem de \emph{declaração} dos membros
--- trocar a ordem dos campos na classe muda a semântica da comparação sem que
nenhuma linha do \texttt{operator<=>} mude.

O valor devolvido não é um \texttt{int}: é uma \textbf{categoria de ordenação},
e a escolha entre as três (\cref{tab:ordenacao}) é uma afirmação sobre o tipo.

\begin{table}[htbp]
\centering
\caption{As três categorias de ordenação e o que cada uma promete.}
\label{tab:ordenacao}
\footnotesize
\begin{tabular}{@{}lcc@{}}
\toprule
\textbf{Categoria} & \textbf{Ordem total} & \textbf{Equivalente = igual} \\
\midrule
\texttt{strong\_ordering}  & sim & \textbf{sim} \\
\texttt{weak\_ordering}    & sim & não \\
\texttt{partial\_ordering} & \textbf{não} & não \\
\bottomrule
\end{tabular}
\end{table}

Os dois casos-limite medidos pelo subprojeto explicam a tabela. Numa comparação
de \emph{strings} sem distinguir maiúsculas, \texttt{"ABC"} e \texttt{"abc"} são
\emph{equivalentes} mas não \emph{iguais} --- isso é \texttt{weak\_ordering}. E
com ponto flutuante, \texttt{1.0 < NaN}, \texttt{1.0 > NaN} e \texttt{1.0 == NaN}
são \textbf{os três falsos} --- é o que \texttt{partial\_ordering} significa, e é
por isso que ordenar um \texttt{vector<double>} com \texttt{NaN} dentro é
\gls{ub}. Prefira a categoria mais forte que for verdadeira: quanto mais forte,
mais os containers podem otimizar.

\begin{regra}[escreveu \texttt{<=>} à mão, escreva o \texttt{==} também]
A ordenação precisa comparar campo a campo; a igualdade quase sempre tem um
atalho barato --- comparar o tamanho primeiro, por exemplo. Por isso o padrão
deriva \texttt{==} de um \texttt{<=>} apenas quando ele é \texttt{= default}.
Um \texttt{<=>} escrito à mão sem \texttt{==} deixa o tipo comparável por ordem
e não comparável por igualdade.
\end{regra}

\paragraph{Candidatos reescritos e revertidos.} O C++20 deriva as comparações
que faltam invertendo os operandos. Declarando apenas duas funções, o compilador
gera oito operações:

\begin{lstlisting}[style=bt, language=C++,
    caption={Comparação heterogênea, em
             \texttt{src/operators/demo\_comparacao.hpp}.}]
struct Metros {
    double valor;

    // Comparação HOMOGÊNEA
    auto operator<=>(const Metros&) const = default;

    // Comparação HETEROGÊNEA: só declaramos Metros <=> double.
    // O compilador deriva `double < Metros` invertendo (candidato revertido).
    std::partial_ordering operator<=>(double d) const { return valor <=> d; }
    bool operator==(double d) const { return valor == d; }
};
\end{lstlisting}

\texttt{m < 7.0} é \emph{reescrito} como \texttt{(m <=> 7.0) < 0}; \texttt{7.0 >
m} é \emph{revertido} para \texttt{0 < (m <=> 7.0)}. Antes do C++20, o mesmo
resultado exigia doze funções escritas à mão.

\subsubsection{Conversões e \texttt{explicit}}

Um construtor de um único argumento \emph{já é} uma conversão implícita, esteja
isso na intenção ou não:

\begin{lstlisting}[style=bt, language=C++,
    caption={Uma palavra separa o tipo forte do tipo decorativo.}]
Metros(double v);            // implícito: qualquer double vira Metros
explicit Metros(double v);   // explícito: só se alguém pedir
\end{lstlisting}

Sem \texttt{explicit}, uma função que espera metros aceita alegremente um valor
em segundos, e compila. É a mesma razão pela qual
\texttt{std::vector<int> v(10)} cria dez zeros e \texttt{std::vector<int> v = 10}
não compila.

A linguagem limita o encadeamento a \textbf{uma} conversão definida pelo
usuário: dado \texttt{struct A \{ A(int); \};} e \texttt{struct B \{ B(const
A\&); \};}, a chamada \texttt{recebe\_b(A\{1\})} compila e
\texttt{recebe\_b(1)} não, porque exigiria duas. Conversões \emph{padrão} não
contam para esse limite --- \texttt{int} para \texttt{double}, derivado para
base, \texttt{T*} para \texttt{const T*}, \emph{enum} para \texttt{int}.

\paragraph{\classe{explicit operator bool}.} É a exceção que confirma a regra:
o único operador de conversão que se marca \texttt{explicit} \textbf{e} ainda
funciona implicitamente, nos contextos ditos \emph{contextualmente
convertíveis} --- \texttt{if}, \texttt{while}, \texttt{for}, \texttt{!},
\texttt{\&\&}, \texttt{||}, \texttt{?:}, \texttt{static\_assert}.

\begin{lstlisting}[style=bt, language=C++,
    caption={O que \texttt{explicit} impede sem impedir o uso pretendido.}]
if (c)      // funciona
int n = c;  // NÃO COMPILA
c + 1;      // NÃO COMPILA
\end{lstlisting}

\begin{porque}[por que os \emph{streams} adotaram \texttt{explicit operator bool}]
Até o C++03, os \emph{streams} convertiam para \texttt{void*} para permitir
\texttt{if (arquivo)}. A consequência era que \texttt{std::cout << std::cin}
compilava --- imprimia um endereço. Sem \texttt{explicit}, um tipo
\texttt{Conexao} com \texttt{operator bool} faria \texttt{conexao + 41} valer 42,
e \texttt{conexaoA == conexaoB} compararia dois \texttt{bool}: duas conexões
fechadas seriam ``iguais''.
\end{porque}

\subsubsection{\texttt{operator()}, \texttt{operator[]} e o \emph{drill-down} de \texttt{operator->}}

Um \emph{functor} --- objeto com \texttt{operator()} --- ganha do ponteiro para
função em tudo o que importa, e a razão é uma só: ele tem \textbf{tipo próprio},
então o compilador sabe exatamente qual código será chamado e pode fazer
\emph{inline}. Um \emph{lambda} \emph{é} um \emph{functor}: o compilador gera
uma classe anônima cujo \texttt{operator()} é \texttt{const} por padrão, e as
capturas viram membros (\cref{sec:lambdas}).

\begin{table}[htbp]
\centering
\caption{As quatro formas de ``algo chamável''.}
\label{tab:chamaveis}
\footnotesize
\begin{tabular}{@{}lccc@{}}
\toprule
\textbf{Forma} & \textbf{Tem estado} & \textbf{Inlinável} & \textbf{Tipo próprio} \\
\midrule
ponteiro para função & não & \textbf{não} & não \\
\emph{functor}       & sim & \textbf{sim} & sim \\
\emph{lambda}        & sim & \textbf{sim} & sim (anônimo) \\
\texttt{std::function} & sim & não & apagado \\
\bottomrule
\end{tabular}
\end{table}

O \texttt{operator[]} multidimensional do C++23 fecha um buraco antigo: até
então, uma matriz tinha três saídas e todas com defeito --- \texttt{m(i,j)}
parece chamada de função, \texttt{m[i][j]} exige uma classe \emph{proxy} por
dimensão, e \texttt{m[\{i,j\}]} constrói um par a cada acesso. Agora
\texttt{vol[i, j, k]} é legítimo.

\begin{armadilha}[\texttt{a[i, j]} mudou de significado entre C++20 e C++23]
Em C++20, \texttt{a[i, j]} era o operador \emph{vírgula} dentro dos colchetes:
avaliava \texttt{i}, descartava, e indexava por \texttt{j}. O C++20 deprecou a
forma e o C++23 reciclou a sintaxe --- de modo que código antigo muda de
comportamento em silêncio ao ser recompilado. O aviso é
\texttt{-Wcomma-subscript}.
\end{armadilha}

O \texttt{operator->} tem uma regra que não existe em nenhum outro ponto da
linguagem: se o que ele devolve \textbf{não} é um ponteiro cru, o compilador
chama \texttt{operator->} de novo no resultado, recursivamente, até chegar a um
ponteiro. É o mecanismo por trás de \texttt{unique\_ptr},
\texttt{optional}, dos iteradores e de qualquer \emph{proxy} de acesso.

\subsubsection{Literais definidos pelo usuário}

Os sufixos de literal permitem que uma grandeza chegue ao código com a unidade
embutida no tipo. O caso mais interessante do subprojeto usa parâmetros de
\emph{template} inteiros para representar as dimensões físicas:

\begin{lstlisting}[style=bt, language=C++,
    caption={Dimensões como expoentes no tipo.}]
template <int Comprimento, int Massa, int Tempo> class Grandeza;
using Velocidade = Grandeza<1, 0, -1>;    // comprimento / tempo
\end{lstlisting}

\begin{lstlisting}[style=bt, language=bash,
    caption={Saída de \texttt{src/operators/}: o tipo do resultado é derivado
             pelo compilador, e o erro dimensional é de compilação.}]
100_km / 2_h      = 13.89 m/s   ← tipo Velocidade, DERIVADO
70_kg * 9.81 m/s² = 686.70 N    ← tipo Forca, derivado
5_km + 300_m      = 5300 m      ✓
5_km + 2_s        → NÃO COMPILA ✗ dimensões incompatíveis
\end{lstlisting}

E o custo é zero: \texttt{sizeof(Comprimento) == sizeof(double)}. É assim que
funcionam as bibliotecas de unidades --- e o próprio \texttt{std::chrono}, que é
um sistema de unidades de tempo com verificação dimensional.

\begin{regra}[todo sufixo de usuário começa com \texttt{\_}]
Sufixos sem \emph{underscore} inicial são reservados para a biblioteca padrão.
\texttt{operator""km} compila hoje e pode colidir com uma adição futura do
padrão amanhã --- e a colisão aparece como um erro incompreensível num arquivo
que ninguém alterou.
\end{regra}

% =============================================================================
\subsection{Ponteiros inteligentes e RAII}
\label{sec:smart-pointers}
% =============================================================================

O subprojeto \projeto{smart-pointers} é \gls{raii} aplicado, e só faz sentido
depois de destrutores e \emph{move} (\cref{sec:value-semantics}). Antes de
mostrar as ferramentas, convém enunciar com precisão o problema que elas
resolvem --- porque quem só viu a solução tende a usá-la como sinônimo de
``ponteiro moderno'', e a errar por excesso.

\subsubsection{O problema, enunciado com precisão}

Chame de \textbf{recurso} qualquer coisa que o programa \emph{adquire} e depois
precisa \emph{devolver}: um bloco de memória, um arquivo aberto, um
\emph{socket}, um \emph{mutex} travado, uma conexão de banco, uma textura na
placa de vídeo. E chame de \textbf{posse} (\emph{ownership}) a resposta à
pergunta ``quem é o responsável por devolver isto, e quando?''.

\begin{porque}[por que a posse precisa estar no tipo, e não no comentário]
Em C, a posse é uma convenção documentada: a página de manual do
\texttt{strdup} diz que quem chama deve dar \texttt{free}. A convenção não é
verificável --- o compilador não a lê, o \emph{linker} não a checa e o teste
não a exercita, porque um vazamento não muda a resposta do programa. É por isso
que os quatro erros da \cref{tab:erros-posse} atravessaram trinta anos de
código: todos eles \emph{compilam}, e três dos quatro passam pelo teste.
\end{porque}

A memória do programa se divide, para este assunto, em duas regiões. A
\textbf{pilha} (\emph{stack}) guarda as variáveis locais; o compilador sabe
onde cada uma nasce e morre, e emite a destruição sozinho. O \textbf{monte}
(\emph{heap}) guarda o que é criado com \texttt{new} e cujo tempo de vida o
compilador \emph{não} consegue deduzir --- e é exatamente aí que a
responsabilidade volta para o programador, com os resultados da
\cref{tab:erros-posse}.

\begin{table}[htbp]
\centering
\caption{Os quatro modos de errar a posse, e por que nenhum deles aparece no
         teste.}
\label{tab:erros-posse}
\footnotesize
\begin{tabular}{@{}lll@{}}
\toprule
\textbf{Erro} & \textbf{O que acontece} & \textbf{Sintoma} \\
\midrule
vazamento (\emph{leak})     & ninguém devolve     & nenhum --- só a memória cresce \\
\emph{double free}          & dois devolvem       & aborto, às vezes \\
\emph{use-after-free}       & usa depois de devolver & lixo, às vezes \\
\emph{dangling pointer}     & aponta para o nada  & \gls{ub} silencioso \\
\bottomrule
\end{tabular}
\end{table}

O argumento inteiro cabe num par de funções:

\begin{lstlisting}[style=bt, language=C++,
    caption={O mesmo recurso, com e sem RAII (\texttt{src/smart-pointers/README.md}).}]
void ruim() {
    Widget* w = new Widget();
    if (condicao) return;        // LEAK
    throw SomeException();       // LEAK
    delete w;                    // só chega aqui no caso feliz
}

void bom() {
    auto w = std::make_unique<Widget>();
    if (condicao) return;        // destrutor chamado
    throw SomeException();       // destrutor chamado
}                                // destrutor chamado
\end{lstlisting}

O que muda não é a quantidade de código: é que na segunda versão \textbf{não
existe caminho de saída que esqueça a liberação}, porque a liberação deixou de
ser um comando e passou a ser uma consequência do escopo. Como o destrutor roda
também durante o desenrolamento da pilha (\cref{sec:errors}), o mesmo mecanismo
cobre o caminho da exceção sem uma linha a mais.

\begin{regra}[a definição operacional de RAII]
\emph{Resource Acquisition Is Initialization}: o construtor adquire o recurso e
o destrutor o devolve, de modo que \textbf{um objeto vivo é um recurso válido}.
A consequência prática é a que interessa: como o compilador garante que o
destrutor de todo objeto automático roda --- em \texttt{return}, em
\texttt{break}, em \texttt{throw} --- ele passa a garantir também a devolução.
O nome é ruim (o que importa é a \emph{destruição}), mas o mecanismo é o
alicerce de metade deste livro.
\end{regra}

\subsubsection{Os três ponteiros}

\begin{table}[htbp]
\centering
\caption{Qual ponteiro para qual relação de posse.}
\label{tab:smart-pointers}
\footnotesize
\begin{tabular}{@{}lll@{}}
\toprule
\textbf{Situação} & \textbf{Use} & \textbf{Custo} \\
\midrule
único dono do objeto         & \classe{std::unique\_ptr} & zero sobre o ponteiro cru \\
posse compartilhada          & \classe{std::shared\_ptr} & bloco de controle + contagem atômica \\
observar sem possuir         & \classe{std::weak\_ptr}   & idem, sem manter vivo \\
acesso sem posse             & \texttt{T*} ou \texttt{T\&} & zero \\
\bottomrule
\end{tabular}
\end{table}

\paragraph{\texttt{unique\_ptr}.} Não é copiável --- a tentativa é erro de
compilação, não bug em produção --- e a transferência de posse é explícita, via
\texttt{std::move}. Depois do \emph{move}, o ponteiro de origem vale
\texttt{nullptr}; usá-lo é \gls{ub}. O \emph{deleter} customizado é o que o
torna um invólucro genérico de recurso, e não apenas de memória:

\begin{lstlisting}[style=bt, language=C++,
    caption={\texttt{unique\_ptr} como guarda de qualquer recurso
             (\texttt{demo\_unique.hpp}, seção 1.5).}]
auto fechar = [](FILE* f){ if (f) std::fclose(f); };
std::unique_ptr<FILE, decltype(fechar)> arq(std::fopen("/dev/null","w"), fechar);
\end{lstlisting}

Duas operações do \texttt{unique\_ptr} parecem sinônimas e não são
(\cref{tab:release-reset}). A distinção é a fronteira entre ``continuo
protegido'' e ``voltei ao C''.

\begin{table}[htbp]
\centering
\caption{\texttt{release()} contra \texttt{reset()} --- a diferença de uma
         chamada de destrutor.}
\label{tab:release-reset}
\footnotesize
\begin{tabular}{@{}lccl@{}}
\toprule
\textbf{Chamada} & \textbf{Destrói?} & \textbf{Fica nulo?} & \textbf{Depois disso} \\
\midrule
\texttt{p.release()} & \textbf{não} & sim & \emph{você} deve dar \texttt{delete} \\
\texttt{p.reset()}   & sim          & sim & nada a fazer \\
\texttt{p.reset(q)}  & sim (o antigo) & não & passa a possuir \texttt{q} \\
\texttt{p.get()}     & não          & não & empréstimo; não guarde \\
\bottomrule
\end{tabular}
\end{table}

\begin{armadilha}[\texttt{release()} é o botão que desliga a proteção]
A linha \texttt{Recurso* raw = p.release();} do \texttt{demo\_unique.hpp} é
seguida, no subprojeto, de um \texttt{delete raw;} explícito e de um comentário
dizendo ``obrigatório --- senão vaza''. Isso não é estilo: \texttt{release()}
existe para entregar o ponteiro a uma API C que vai assumir a posse, e em
qualquer outro uso ele apenas reintroduz os quatro erros da
\cref{tab:erros-posse}. Quando a intenção é só \emph{emprestar}, o certo é
\texttt{get()} --- ou, melhor ainda, passar \texttt{*p} por referência.
\end{armadilha}

Existe também \texttt{std::unique\_ptr<T[]>}, que chama \texttt{delete[]} em vez
de \texttt{delete} --- e essa é a única razão para ele existir, porque trocar um
pelo outro é \gls{ub}. Ainda assim, ele quase nunca é a resposta certa: não
guarda o tamanho, não tem \texttt{begin()}/\texttt{end()} e não cresce.
\classe{std::vector} faz tudo isso pelo mesmo preço.

\paragraph{\texttt{shared\_ptr}.} Contagem de referências: o objeto morre quando
o último \texttt{shared\_ptr} morre. O que ele carrega além do ponteiro é o
\textbf{bloco de controle}, um objeto separado com duas contagens --- quantos
\emph{shared} e quantos \emph{weak} apontam para o alvo --- e o \emph{deleter}.
\texttt{make\_shared} funde esse bloco com o objeto numa alocação só, contra as
duas de \texttt{std::shared\_ptr<T> p(new T)}, e é também seguro em face de
exceções.

\paragraph{\texttt{weak\_ptr}.} Observa sem possuir, e existe por causa de um
problema que a contagem de referências não resolve sozinha: o \textbf{ciclo}.

\begin{lstlisting}[style=bt, language=C++,
    caption={O ciclo, e a mudança de uma palavra que o quebra
             (\texttt{demo\_weak.hpp}, seções 3.2 e 3.3).}]
struct No { std::shared_ptr<No> proximo; };   // a → b → a: use_count nunca zera
struct No { std::weak_ptr<No>   proximo; };   // não incrementa a contagem
\end{lstlisting}

O subprojeto roda os dois casos em sequência, e a saída é a prova: no caso do
ciclo os destrutores \emph{não} aparecem.

\begin{lstlisting}[style=bt, language=bash,
    caption={Saída de \texttt{dist/bin/smart\_pointers}, seções 3.2 e 3.3.
             Note quais destrutores faltam.}]
── 3.2 Ciclo com shared_ptr → MEMORY LEAK ──
  [NodoComCiclo 'A-ciclo'] construído
  [NodoComCiclo 'B-ciclo'] construído
  use_count(a) = 2
  use_count(b) = 2
                              ← nenhum destrutor: os dois vazaram

── 3.3 Ciclo com weak_ptr → SEM LEAK ──
  [NodoSemCiclo 'A-ok'] construído
  [NodoSemCiclo 'B-ok'] construído
  use_count(a) = 1
  use_count(b) = 1
  a->proximo = 'B-ok'
  [NodoSemCiclo 'B-ok'] destruído
  [NodoSemCiclo 'A-ok'] destruído
\end{lstlisting}

Para usar o objeto é preciso \texttt{lock()}, que devolve um
\texttt{shared\_ptr} vazio se ele já morreu --- e é essa obrigação de verificar
que torna o \texttt{weak\_ptr} seguro. \texttt{expired()} responde a mesma
pergunta sem produzir o \texttt{shared\_ptr}, e por isso \textbf{não} serve como
teste antes do uso: entre o \texttt{expired()} e a desreferência, outra
\emph{thread} pode ter soltado a última referência. O idioma correto é um só,
e já traz o resultado consigo:

\begin{lstlisting}[style=bt, language=C++,
    caption={O único uso seguro de um \texttt{weak\_ptr}.}]
if (auto p = w.lock()) { p->usar(); }   // p mantém vivo enquanto durar o if
\end{lstlisting}

\begin{armadilha}[\texttt{use\_count()} não conta os \texttt{weak\_ptr}]
Criar um \texttt{weak\_ptr} a partir de um \texttt{shared\_ptr} não altera
\texttt{use\_count()} --- ele continua respondendo 1, como imprime a seção 3.1
do subprojeto. Isso é intencional (o \texttt{weak\_ptr} não mantém o objeto
vivo), mas contraria a leitura ingênua do nome. O que os \emph{weak} mantêm
vivo é o \emph{bloco de controle}, não o objeto --- consequência que a
\cref{sec:sp-medido} mede em bytes.
\end{armadilha}

\subsubsection{O custo, medido}
\label{sec:sp-medido}

\begin{aviso}
Todos os números desta seção foram medidos neste repositório, em
\texttt{AMD Ryzen 9 5900X}, GCC 13.3.0, libstdc++ 13, glibc 2.39,
\texttt{-O2 -std=c++20}. Números de desempenho não são propriedades da
linguagem: são propriedades de uma combinação de compilador, biblioteca,
máquina e carga. O que se transporta de uma medição para outra é a
\emph{ordem de grandeza} e, sobretudo, o \emph{método}.
\end{aviso}

A afirmação ``\texttt{unique\_ptr} tem custo zero'' é repetida em todo lugar e
quase nunca demonstrada. Demonstrar é fácil: basta comparar o código de máquina
das duas versões da mesma função.

\begin{lstlisting}[style=bt, language=C++,
    caption={As duas funções comparadas.}]
struct W { int a; W(int x):a(x){} ~W(); };
int cru()   { W* p = new W(7); int v = p->a; delete p; return v; }
int smart() { auto p = std::make_unique<W>(7); return p->a; }
\end{lstlisting}

\begin{lstlisting}[style=bt, language=bash,
    caption={\texttt{g++ -O2 -S}: a sequência de instruções das duas funções é
             idêntica, instrução por instrução --- 14 cada.}]
_Z3cruv:                 _Z5smartv:
  endbr64                  endbr64
  pushq  %rbx              pushq  %rbx
  movl   $4, %edi          movl   $4, %edi
  call   _Znwm@PLT         call   _Znwm@PLT      ← operator new
  movl   $7, (%rax)        movl   $7, (%rax)
  movq   %rax, %rbx        movq   %rax, %rbx
  movq   %rax, %rdi        movq   %rax, %rdi
  call   _ZN1WD1Ev@PLT     call   _ZN1WD1Ev@PLT  ← ~W()
  movq   %rbx, %rdi        movq   %rbx, %rdi
  movl   $4, %esi          movl   $4, %esi
  call   _ZdlPvm@PLT       call   _ZdlPvm@PLT    ← operator delete
  movl   $7, %eax          movl   $7, %eax
  popq   %rbx              popq   %rbx
  ret                      ret
\end{lstlisting}

Não é ``quase igual'': é a mesma sequência. O \texttt{unique\_ptr} desaparece
por completo na compilação, e o que sobra é o \texttt{new}/\texttt{delete} que
você teria escrito à mão --- com a diferença de que o \texttt{delete} agora é
emitido pelo compilador em \emph{todos} os caminhos de saída, inclusive nos que
você esqueceria.

O tamanho segue a mesma lógica: o que não é usado não é pago
(\cref{tab:sp-sizeof}).

\begin{table}[htbp]
\centering
\caption{\texttt{sizeof} medido com \texttt{g++ -std=c++20} e libstdc++ 13.}
\label{tab:sp-sizeof}
\footnotesize
\begin{tabular}{@{}lrl@{}}
\toprule
\textbf{Tipo} & \textbf{bytes} & \textbf{Por quê} \\
\midrule
\texttt{int*}                              &  8 & referência \\
\texttt{std::unique\_ptr<W>}               &  8 & \emph{deleter} sem estado, elidido \\
\texttt{std::unique\_ptr<W[]>}             &  8 & idem \\
\texttt{std::unique\_ptr<W, lambda>}       &  8 & \emph{lambda} sem captura é vazia \\
\texttt{std::unique\_ptr<W, void(*)(W*)>}  & \textbf{16} & o ponteiro de função \emph{tem} estado \\
\texttt{std::shared\_ptr<W>}               & 16 & alvo + bloco de controle \\
\texttt{std::weak\_ptr<W>}                 & 16 & idem \\
\bottomrule
\end{tabular}
\end{table}

\begin{armadilha}[o \emph{deleter} pode dobrar o tamanho do ponteiro]
As duas formas de guardar um \texttt{FILE*} parecem equivalentes e não são.
Com uma \emph{lambda} sem captura, o \emph{deleter} é um tipo vazio e a
\emph{empty base optimization} o faz sumir: 8 bytes. Com
\texttt{decltype(\&std::fclose)}, o \emph{deleter} é um ponteiro de função ---
um dado real, guardado em cada instância: 16 bytes, o dobro, e um campo a mais
para copiar em cada \emph{move}. Num \texttt{std::vector} de um milhão desses
guardas, a diferença é 8~MB de nada.
\end{armadilha}

O \texttt{shared\_ptr}, esse, custa. A \cref{tab:sp-aloc} conta as chamadas a
\texttt{operator new} substituindo-o globalmente por uma versão que incrementa
um contador --- técnica que reaparece no \cref{sec:memory}.

\begin{table}[htbp]
\centering
\caption{Alocações e bytes por criação, medidos com \texttt{operator new}
         instrumentado sobre um objeto de 64~bytes.}
\label{tab:sp-aloc}
\footnotesize
\begin{tabular}{@{}lrrl@{}}
\toprule
\textbf{Expressão} & \textbf{Alocações} & \textbf{Bytes} & \textbf{Sobrecusto} \\
\midrule
\texttt{std::make\_unique<Big>()}          & 1 & 64 & nenhum \\
\texttt{std::make\_shared<Big>()}          & 1 & 80 & +16 (bloco embutido) \\
\texttt{std::shared\_ptr<Big>(new Big())}  & \textbf{2} & 88 & +24 e uma segunda \texttt{malloc} \\
\bottomrule
\end{tabular}
\end{table}

E a contagem custa tempo. Vinte milhões de iterações de cada laço, tempo médio
por operação:

\begin{lstlisting}[style=bt, language=bash,
    caption={Custo por operação, 20 milhões de iterações, \texttt{-O2}.}]
copia de shared_ptr (incr + decr atomico) : 4,43 ns/op
weak_ptr::lock()                          : 5,27 ns/op
copia de ponteiro cru                     : 0,00 ns/op   ← ver a armadilha abaixo
\end{lstlisting}

\begin{armadilha}[o ``0,00 ns'' não é velocidade: é ausência]
O laço do ponteiro cru mediu zero porque o compilador o \emph{apagou}: copiar um
ponteiro e somar um campo que ninguém observa não tem efeito colateral, e
\texttt{-O2} remove o laço inteiro. O laço do \texttt{shared\_ptr} sobreviveu
justamente porque a contagem atômica \textbf{é} um efeito observável, que o
compilador não tem licença para eliminar. Toda medição de microbenchmark tem
esse risco, e o sintoma é sempre o mesmo: um número absurdamente bom. O hábito
que evita a conclusão errada é conferir o assembly antes de acreditar --- ou
usar \texttt{benchmark::DoNotOptimize}. A conclusão honesta desta medição é
``copiar \texttt{shared\_ptr} custa cerca de 4~ns e copiar ponteiro cru custa
menos que a resolução do experimento'', não ``é infinitamente mais rápido''.
\end{armadilha}

\subsubsection{A medição que contrariou o conselho: \texttt{make\_shared} e o \texttt{weak\_ptr}}

Todo material sobre \texttt{shared\_ptr} --- inclusive o
\texttt{README.md} deste subprojeto --- diz para preferir
\texttt{make\_shared}, e a \cref{tab:sp-aloc} confirma a razão: uma alocação em
vez de duas. Mas a fusão do bloco de controle com o objeto tem uma consequência
que o conselho não menciona, e que só aparece quando se \emph{mede a memória
liberada} em vez de contar as alocações.

O experimento: um objeto de 1~KB, um \texttt{weak\_ptr} que sobrevive ao último
\texttt{shared\_ptr}, e \texttt{operator delete} instrumentado para contar as
liberações.

\begin{lstlisting}[style=bt, language=bash,
    caption={Liberações contadas com \texttt{operator delete} instrumentado,
             objeto de 1024 bytes, um \texttt{weak\_ptr} vivo.}]
[make_shared]      dentro do escopo:            alocacoes=1  liberacoes=0
[make_shared]      ~Grande() rodou
[make_shared]      apos o ultimo shared morrer: liberacoes=0   ← 1 KB retido!
[make_shared]      apos o weak_ptr morrer:      liberacoes=1

[shared_ptr(new)]  dentro do escopo:            alocacoes=2
[shared_ptr(new)]  ~Grande() rodou
[shared_ptr(new)]  apos o ultimo shared morrer: liberacoes=1   ← 1 KB devolvido
[shared_ptr(new)]  apos o weak_ptr morrer:      liberacoes=2
\end{lstlisting}

Os dois destroem o objeto na hora certa --- \texttt{\~{}Grande()} roda nos dois
casos assim que o último \texttt{shared\_ptr} morre. O que muda é a
\emph{memória}: com \texttt{make\_shared} o bloco é um só, e enquanto houver um
\texttt{weak\_ptr} vivo o alocador não pode devolver nada, nem o quilobyte do
objeto já destruído. Com \texttt{shared\_ptr(new)}, o objeto tem alocação
própria e volta imediatamente; só os 24~bytes do bloco de controle esperam pelo
\emph{weak}.

\begin{regra}[\texttt{make\_shared} por padrão, \texttt{shared\_ptr(new)} quando o \emph{weak} sobrevive]
Prefira \texttt{make\_shared}: uma alocação em vez de duas, melhor localidade e
segurança em face de exceções. \textbf{Exceto} quando o objeto for grande e
houver \texttt{weak\_ptr} com expectativa de vida maior que a dele --- caches,
tabelas de observadores, grafos com arestas fracas. Nesse caso a alocação extra
compra a devolução imediata do objeto, e é o negócio certo. Este é o formato
típico de um conselho bom: verdadeiro no caso comum, falso num caso nomeável, e
decidível por medição.
\end{regra}

\subsubsection{A contagem é atômica --- e não dá mais para desligá-la}

Durante anos valeu o seguinte folclore: a libstdc++ detecta em tempo de execução
se o programa é \emph{single-threaded} e, nesse caso, usa incrementos comuns em
vez de instruções atômicas; logo, quem não liga \texttt{-pthread} não paga pela
atomicidade. Isso era verdade, e deixou de ser. A verificação era
\texttt{\_\_gthread\_active\_p()}, que perguntava se o símbolo
\texttt{pthread\_create} havia sido resolvido; desde a glibc~2.34 a
\texttt{libpthread} foi fundida na \texttt{libc}, e o símbolo está sempre lá.

\begin{lstlisting}[style=bt, language=bash,
    caption={Contagem de prefixos \texttt{lock} numa cópia de
             \texttt{shared\_ptr}, glibc 2.39.}]
$ g++ -O2 -S copia.cpp          && grep -c 'lock ' copia.s
3
$ g++ -O2 -pthread -S copia.cpp && grep -c 'lock ' copia.s
3
\end{lstlisting}

Os mesmos três \texttt{lock xaddl} nos dois casos. A consequência prática: os
4,43~ns por cópia da seção anterior são o preço real, mesmo em programa de uma
\emph{thread} só. E a consequência didática é mais importante que a técnica ---
\textbf{um conselho de desempenho tem prazo de validade}, e o único modo de
saber se o dele venceu é remedir.

\begin{regra}[passe \texttt{shared\_ptr} por \texttt{const\&} quando não for compartilhar]
Cada passagem por valor é um par de operações atômicas: incremento na entrada,
decremento na saída. Numa função chamada em laço, isso é o custo dominante.
Passe por valor \emph{só} quando a função for de fato guardar uma referência
nova; para apenas usar o objeto, o parâmetro certo é \texttt{const T\&} ---
ou \texttt{T\&}, que nem menciona posse.
\end{regra}

\subsubsection{\texttt{enable\_shared\_from\_this} e o construtor de \emph{aliasing}}

Duas facilidades do \texttt{shared\_ptr} que o subprojeto cobre e que costumam
faltar em material introdutório --- ambas resolvem problemas reais, e ambas têm
uma armadilha de mão dupla.

\paragraph{Um objeto que precisa de um \texttt{shared\_ptr} de si mesmo.} O
caso aparece sempre que um método precisa se registrar em algum lugar
(\emph{callback} assíncrono, lista de observadores). A tentação é
\texttt{return std::shared\_ptr<T>(this)} --- e isso cria um \textbf{segundo
bloco de controle} para o mesmo endereço. Dois blocos, duas contagens, duas
chegadas a zero: \emph{double free}. A solução é herdar de
\texttt{std::enable\_shared\_from\_this<T>}, que guarda um \texttt{weak\_ptr}
para o bloco existente e o promove sob demanda:

\begin{lstlisting}[style=bt, language=C++,
    caption={\texttt{shared\_from\_this} devolve uma referência ao \emph{mesmo}
             bloco (\texttt{demo\_shared.hpp}, seção 2.5).}]
struct Observavel : std::enable_shared_from_this<Observavel> {
    std::shared_ptr<Observavel> get_self() { return shared_from_this(); }
};
// medido: use_count 1 → 2 dentro do escopo de get_self() → 1 de volta
\end{lstlisting}

\begin{armadilha}[\texttt{shared\_from\_this} num objeto que não é de um \texttt{shared\_ptr}]
O mecanismo só funciona se o objeto \emph{já} for gerenciado por um
\texttt{shared\_ptr} --- é dele que vem o bloco de controle. Chamar
\texttt{shared\_from\_this()} numa variável de pilha, ou dentro do próprio
construtor (quando o \texttt{shared\_ptr} externo ainda não existe), lança
\texttt{std::bad\_weak\_ptr}. É por isso que classes assim costumam ter
construtor privado e uma \emph{factory} estática que devolve
\texttt{make\_shared}: a fábrica torna a pré-condição impossível de violar.
\end{armadilha}

\paragraph{Um ponteiro para dentro do objeto.} O \emph{alias constructor}
\texttt{std::shared\_ptr<U>(dono, \&dono->membro)} produz um
\texttt{shared\_ptr} que \emph{aponta} para o membro mas \emph{conta} no bloco
do objeto inteiro. Serve para entregar uma parte sem entregar o todo:

\begin{lstlisting}[style=bt, language=bash,
    caption={Saída medida de \texttt{dist/bin/smart\_pointers}, seção 2.6.}]
  [Carro 'Fusca'] construído
  carro.use_count() = 2 (carro + motor_ptr)
  motor_ptr->rpm = 3000
  carro.use_count() após reset: 1
  motor_ptr->rpm ainda acessível: 3000     ← o Carro continua vivo
  [Carro 'Fusca'] destruído                ← só agora, com motor_ptr saindo
\end{lstlisting}

\begin{armadilha}[um \emph{alias} pequeno segura um objeto grande]
Exatamente a propriedade que torna o \emph{alias constructor} seguro --- o
membro nunca sobrevive ao dono --- é a que causa retenção acidental: um
\texttt{shared\_ptr<int>} de 16~bytes apontando para um campo de uma imagem de
40~MB mantém os 40~MB vivos. É o mesmo fenômeno do \texttt{make\_shared} com
\texttt{weak\_ptr} da seção anterior, e o mesmo remédio: saber que existe, e
copiar o valor quando o que se quer é só o valor.
\end{armadilha}

\subsubsection{Como passar um ponteiro inteligente para uma função}

A assinatura de uma função é onde a posse fica \emph{documentada pelo
compilador}. A \cref{tab:sp-passagem} traduz cada forma; ela é a versão para
recursos da \cref{tab:passagem}.

\begin{table}[htbp]
\centering
\caption{O que cada assinatura comunica (\texttt{demo\_patterns.hpp},
         ``regras de ouro'').}
\label{tab:sp-passagem}
\footnotesize
\begin{tabular}{@{}ll@{}}
\toprule
\textbf{Assinatura} & \textbf{Significado} \\
\midrule
\texttt{f(Widget\&)}                 & só usa; posse não é assunto da função \\
\texttt{f(const Widget\&)}           & só lê --- \textbf{o padrão} \\
\texttt{f(Widget*)}                  & só usa, e pode não haver objeto \\
\texttt{f(unique\_ptr<W>)}           & \textbf{toma} a posse exclusiva \\
\texttt{f(unique\_ptr<W>\&)}         & pode \emph{trocar} o ponteiro de quem chamou \\
\texttt{f(shared\_ptr<W>)}           & passa a compartilhar (custa duas atômicas) \\
\texttt{f(const shared\_ptr<W>\&)}   & pode vir a compartilhar; não paga se não vier \\
\bottomrule
\end{tabular}
\end{table}

\begin{regra}[a maioria das funções não deveria ver um ponteiro inteligente]
Se a função apenas \emph{usa} o objeto, ela não tem opinião sobre posse e o
parâmetro certo é \texttt{const Widget\&}. Receber \texttt{const
shared\_ptr<Widget>\&} onde bastava uma referência acopla a função a uma
política de posse que ela não precisa conhecer --- e impede chamá-la com um
objeto de pilha, ou com um membro de outro objeto. Ponteiro inteligente em
assinatura é para quem \emph{vai mexer na posse}.
\end{regra}

\subsubsection{Três padrões que dependem de \texttt{unique\_ptr}}

\paragraph{Guarda genérico.} Qualquer recurso com função de fechamento vira
RAII com um \emph{deleter}: \texttt{FILE*}, \emph{socket}, \emph{mutex},
\emph{handle} do sistema operacional. É o mesmo mecanismo que o
\cref{sec:plugin-registry} usa para \texttt{dlclose}.

\paragraph{Pimpl.} Um \texttt{unique\_ptr} para uma \texttt{struct Impl} apenas
declarada no cabeçalho quebra a dependência de compilação entre a interface e o
que ela usa: mudar os campos de \texttt{Impl} deixa de obrigar a recompilar
quem inclui o cabeçalho. O detalhe que faz o idioma funcionar --- e que é a
causa de metade dos erros de quem o tenta pela primeira vez --- é que o
destrutor precisa ser \emph{declarado} no cabeçalho e \emph{definido} no
\texttt{.cpp}, onde \texttt{Impl} já é um tipo completo. A razão é mecânica: o
destrutor implícito instancia \texttt{default\_delete<Impl>}, que faz
\texttt{sizeof(Impl)} para verificar que o tipo não é incompleto.

\begin{lstlisting}[style=bt, language=bash,
    caption={Erro real do GCC 13.3 ao esquecer o destrutor fora de linha ---
             note que ele aponta para dentro do \texttt{<memory>}, não para o
             seu código.}]
/usr/include/c++/13/bits/unique_ptr.h:97:23: error:
    invalid application of 'sizeof' to incomplete type 'Widget::Impl'
\end{lstlisting}

\begin{armadilha}[o Pimpl que ``compila em um \texttt{.cpp} e quebra no outro'']
Enquanto o único \texttt{.cpp} que destrói o objeto for aquele em que
\texttt{Impl} está definido, tudo funciona. O erro aparece no primeiro
\emph{outro} arquivo que deixe um \texttt{Widget} sair de escopo --- muitas
vezes semanas depois, num arquivo que ninguém tocou. A profilaxia é declarar os
cinco (\cref{sec:value-semantics}) no cabeçalho e defini-los, mesmo que com
\texttt{= default}, no \texttt{.cpp}.
\end{armadilha}

\paragraph{\emph{Factory}.} Devolver \texttt{std::unique\_ptr<Base>} torna a
posse inequívoca no tipo de retorno: quem chamou é dono, e o compilador não
deixa esquecer. É a assinatura que aparece em \cref{sec:patterns} e
\cref{sec:plugin-registry}.

\begin{regra}[\texttt{new} e \texttt{delete} crus não aparecem em código de aplicação]
Eles pertencem ao interior de um tipo que gerencia recurso --- e, na prática, ao
interior da biblioteca padrão. Código de aplicação usa \texttt{make\_unique},
\texttt{make\_shared} ou um container. O \cref{sec:memory} mostra o que acontece
dentro de \texttt{new}, e é a exceção que justifica a regra.
\end{regra}

% =============================================================================
\subsection{Arquivos e \emph{streams}: RAII com estado}
\label{sec:streams-raii}
% =============================================================================

O subprojeto \projeto{file\_io} entra aqui, e não no capítulo da biblioteca
padrão, por um motivo de assunto: um \emph{stream} é o primeiro recurso RAII
que o leitor encontra na prática, e é o único da biblioteca padrão que
\textbf{carrega estado de erro em vez de lançar}. Quem entendeu
\texttt{unique\_ptr} tem metade do modelo; a outra metade --- por que
\texttt{ifs >> x} não lança quando falha, e o que exatamente acontece então ---
é o que esta seção fecha. O catálogo das classes, com a hierarquia desenhada,
está na \cref{sec:file-io}; aqui interessa o mecanismo e o que ele custa.

\begin{porque}[por que um arquivo é mais difícil que um bloco de memória]
Memória tem um dono e uma operação de devolução. Um arquivo tem quatro coisas a
mais, e cada uma é uma fonte de bug distinta. \textbf{(1)}~Um
\emph{descritor} --- um número inteiro que o núcleo do sistema entrega no
\texttt{open} e cobra de volta no \texttt{close}; o processo tem uma cota deles,
e vazá-los derruba o programa muito antes de a memória acabar.
\textbf{(2)}~Um \emph{buffer}: o que você escreveu ainda não chegou ao disco, e
só chega no \emph{flush}. \textbf{(3)}~Um \emph{cursor}, a posição corrente de
leitura e a de escrita, que são duas e independentes. \textbf{(4)}~Um
\emph{estado}: a operação anterior pode ter falhado, e o \emph{stream} guarda
isso em vez de avisar. Ignorar qualquer um dos quatro produz um programa que
funciona no teste e perde dados em produção.
\end{porque}

\subsubsection{O \emph{stream} como objeto RAII}

\texttt{std::ofstream} adquire no construtor e devolve no destrutor, exatamente
como o \texttt{unique\_ptr}: abre o arquivo, e no fim do escopo faz
\emph{flush} do \emph{buffer} e fecha o descritor. É a razão pela qual o
\texttt{demo\_fstream.hpp} abre cada arquivo dentro de um bloco próprio ---
o fechamento acontece na chave, não numa chamada que se possa esquecer.

\begin{lstlisting}[style=bt, language=C++,
    caption={O bloco \emph{é} o \texttt{close} (\texttt{demo\_fstream.hpp},
             seção 1).}]
{
    std::ofstream ofs(path);              // abre
    ofs << "Linha 1: Hello, C++!\n";
}                                         // flush + close, inclusive se lançar
\end{lstlisting}

\begin{armadilha}[o destrutor fecha, mas não conta se deu errado]
O destrutor do \texttt{ofstream} faz \emph{flush} e engole qualquer falha ---
disco cheio, cota estourada, sistema de arquivos remoto que caiu. Como o
destrutor não pode lançar (\cref{sec:errors}), o erro simplesmente não é
reportado, e o programa termina com código de saída zero tendo escrito um
arquivo truncado. Quando a gravação \emph{importa}, chame \texttt{close()}
explicitamente e teste o \emph{stream} depois --- é a única forma de saber.
\end{armadilha}

\subsubsection{Os quatro bits de estado}

Um \emph{stream} não lança por padrão. Ele acende bits, e continua aceitando
operações que, a partir dali, não fazem nada. Os bits são quatro
(\cref{tab:stream-bits}), e a confusão entre \texttt{fail} e \texttt{eof} é a
origem do erro mais comum de todo o assunto.

\begin{table}[htbp]
\centering
\caption{Os bits de estado de um \emph{stream}, e o que cada um significa.}
\label{tab:stream-bits}
\footnotesize
\begin{tabular}{@{}lll@{}}
\toprule
\textbf{Consulta} & \textbf{Aceso quando} & \textbf{Recuperável?} \\
\midrule
\texttt{good()} & nenhum dos outros está aceso & --- \\
\texttt{eof()}  & a última leitura tocou o fim do arquivo & sim \\
\texttt{fail()} & a conversão falhou, ou o arquivo não abriu & sim, com \texttt{clear()} \\
\texttt{bad()}  & erro de integridade do \emph{stream} & \textbf{não} \\
\bottomrule
\end{tabular}
\end{table}

O subprojeto imprime os quatro para um arquivo que não existe, e o resultado
esclarece a diferença entre ``acabou'' e ``deu errado'':

\begin{lstlisting}[style=bt, language=bash,
    caption={Saída medida de \texttt{dist/bin/file\_io}, demo 2, seção 5.}]
── abrindo /tmp/nao_existe.txt ──
  arquivo não abriu (is_open() = false)
  good()=false fail()=true bad()=false eof()=false
\end{lstlisting}

\texttt{fail()} está aceso, \texttt{eof()} não: não se chegou ao fim de coisa
alguma, a abertura é que falhou. O mesmo par aparece ao ler um \texttt{int} de
onde há a palavra \texttt{abc} --- \texttt{fail()=true}, \texttt{eof()=false},
e o \texttt{int} fica \textbf{inalterado}, não zerado. Só depois de
\texttt{clear()} o \emph{stream} volta a funcionar; note, na saída do demo, que
o texto \texttt{abc} continua no \emph{buffer} e é lido em seguida como
\texttt{std::string}.

\begin{armadilha}[\texttt{while (!ifs.eof())} lê a última linha duas vezes]
\texttt{eof()} só acende \emph{depois} de uma leitura que tentou passar do fim.
No laço \texttt{while (!ifs.eof()) \{ std::getline(ifs, l); usa(l); \}}, a
última iteração lê nada, falha, e ainda assim entra no corpo com o valor antigo
de \texttt{l}. A forma correta põe a leitura na condição ---
\texttt{while (std::getline(ifs, l))} --- porque o \emph{stream} converte para
\texttt{bool} testando \texttt{!fail()}, que é a pergunta certa: ``a leitura que
eu acabei de pedir funcionou?''.
\end{armadilha}

Quem prefere exceções pode pedi-las com \texttt{exceptions(std::ios::failbit |
std::ios::badbit)}, e aí a abertura malsucedida lança
\texttt{std::ios\_base::failure}. O subprojeto faz as duas coisas para
comparar --- e a mensagem que sai, \texttt{basic\_ios::clear: iostream error},
ilustra bem por que a via das exceções é pouco usada aqui: ela não diz qual
arquivo, nem por quê. Para saber isso é preciso o \texttt{errno}, ou o
\texttt{std::error\_code} do \cabecalho{filesystem}.

\subsubsection{Bytes não são caracteres}

O demo escreve quatro linhas de texto e pergunta o tamanho ao sistema de
arquivos. O número que volta não é o número de caracteres:

\begin{lstlisting}[style=bt, language=bash,
    caption={\texttt{dist/bin/file\_io}, demo 2, seção 1, confrontado com
             \texttt{wc}.}]
  tamanho final: 89 bytes          ← fs::file_size()
$ wc -m -c /tmp/cpp_academy_texto.txt
     87      89                    ← 87 caracteres, 89 bytes
\end{lstlisting}

A diferença de dois bytes está numa linha só: o travessão de
\texttt{"Linha 2: 42 --- 3.14"}. U+2014 ocupa \textbf{três} bytes em UTF-8, e
todas as APIs de arquivo contam bytes.

\begin{armadilha}[o caractere de três bytes entre aspas simples]
Esta não é uma preocupação acadêmica: foi um bug real deste repositório, em
\texttt{src/http-client/main.cpp}. A linha escrevia
\texttt{std::string(n, '─')} para desenhar uma régua --- e U+2500, como U+2014,
tem três bytes. Entre \emph{aspas simples} isso não é um \texttt{char}: é uma
\emph{multi-character constant}, cujo valor é definido pela implementação e que
o construtor trunca para um byte, produzindo lixo. A mesma função continha o
segundo bug clássico do assunto: \texttt{40 - title.size()}, com
\texttt{size()} devolvendo \texttt{std::size\_t} --- aritmética \textbf{sem
sinal}, que para um título de 41 caracteres não dá $-1$, dá
\texttt{18446744073709551615}. Os dois compilavam sem um \emph{warning} sequer.
Para repetir uma sequência multibyte, o certo é
\texttt{std::string} concatenada em laço, ou uma \texttt{std::string\_view}
com o literal.
\end{armadilha}

\subsubsection{Texto, binário e o cursor}

Em modo texto o \emph{stream} formata: \texttt{ofs << 42} escreve dois
caracteres. Em modo binário ele copia bytes: \texttt{ofs.write(ptr, n)} despeja
a representação de memória tal como está. O demo grava três registros e mede:

\begin{lstlisting}[style=bt, language=bash,
    caption={\texttt{dist/bin/file\_io}, demo 2, seção 4.}]
  3 registros escritos (24 bytes cada, total=72 bytes)
  [ 1] id=  1 valor=3.14 nome='alpha'
  [ 2] id=  2 valor=2.71 nome='beta'
  [ 3] id=  3 valor=1.41 nome='gamma'
  gcount()=0 (bytes lidos na última operação)
  acesso aleatório → 2º registro: id=2 nome='beta'
\end{lstlisting}

Os 24 bytes por registro são \texttt{uint32\_t} + \texttt{float} +
\texttt{char[16]} --- 4 + 4 + 16, sem preenchimento, porque o alinhamento mais
restritivo da estrutura é 4 e o tamanho já é múltiplo dele. Não conte com isso:
é uma coincidência aritmética desta estrutura, não uma regra
(\cref{sec:layout} trata do assunto).

\begin{armadilha}[\texttt{gcount()} depois do laço vale zero]
A saída acima diz \texttt{gcount()=0} logo depois de ter lido três registros de
24 bytes. Não é bug: \texttt{gcount()} devolve os bytes da \emph{última}
operação de entrada não formatada, e a última foi a que encerrou o laço --- a
leitura que encontrou o fim do arquivo e leu zero bytes. Para saber quanto veio
de um \texttt{read} parcial, é preciso consultar \texttt{gcount()}
\textbf{dentro} do laço, imediatamente após cada leitura. É um caso exemplar da
classe de erro deste capítulo: o programa está certo, a saída está certa, e a
leitura ingênua está errada.
\end{armadilha}

Gravar uma \texttt{struct} com \texttt{reinterpret\_cast} congela o
preenchimento, o tamanho dos tipos e a ordem dos bytes (\emph{endianness}) da
máquina que gravou: o arquivo de 72 bytes acima não é legível por um programa
compilado com outro alinhamento, nem por uma máquina \emph{big-endian}. A
\cref{sec:io-streams} volta ao ponto; dentro de uma máquina, para \emph{cache}
de curta duração, a técnica é legítima.

Um \texttt{fstream} aberto com \texttt{ios::in | ios::out} tem dois cursores
independentes --- o de leitura (\texttt{tellg}/\texttt{seekg}) e o de escrita
(\texttt{tellp}/\texttt{seekp}) --- e uma propriedade que costuma surpreender:
escrever em um \emph{offset} \textbf{substitui} os bytes que estavam lá. No
demo, cinco caracteres gravados no início transformam a linha
\texttt{ITEM\_0} em \texttt{MODIF0}: ``MODIF'' cobriu ``ITEM\_'' e o
\texttt{0} sobreviveu. Não existe ``inserir no meio de um arquivo''; inserir é
reescrever o resto --- e é por isso que editores de texto carregam o arquivo
inteiro na memória.

\subsubsection{O que os \emph{streams} custam, medido}

Três conselhos circulam sobre desempenho de \emph{streams}. Medidos nesta
máquina (mesmas condições da \cref{sec:sp-medido}), um se confirma com folga,
um se inverte e um encolhe.

\begin{table}[htbp]
\centering
\caption{200\,000 linhas escritas, lidas ou concatenadas; média de duas
         execuções, \texttt{-O2}, disco local.}
\label{tab:io-bench}
\footnotesize
\begin{tabular}{@{}llrl@{}}
\toprule
\textbf{Operação} & \textbf{Forma} & \textbf{Tempo} & \textbf{Veredito} \\
\midrule
escrever & \texttt{ofs << ... << '\textbackslash n'} &   8,8~ms & --- \\
escrever & \texttt{ofs << ... << std::endl}          & 280~ms   & \textbf{32$\times$ mais lento} \\
\midrule
ler      & \texttt{getline} linha a linha            &   2,9~ms & --- \\
ler      & \texttt{oss << ifs.rdbuf()}               &   2,4~ms & 1,2$\times$ mais rápido \\
\midrule
montar   & \texttt{std::string} com \texttt{+=}      &   3,6~ms & \textbf{mais rápido} \\
montar   & \texttt{std::ostringstream}               &   6,6~ms & 1,8$\times$ mais lento \\
\bottomrule
\end{tabular}
\end{table}

\paragraph{O conselho que se confirmou.} \texttt{std::endl} não é ``quebra de
linha'': é quebra de linha \textbf{mais \emph{flush}}, e cada \emph{flush} é
uma chamada de sistema. Trinta e duas vezes mais lento é a diferença entre
escrever um \emph{log} de 200 mil linhas em 9~ms e em 0,28~s.

\begin{regra}[\texttt{'\textbackslash n'} sempre; \texttt{std::endl} só quando o \emph{flush} for o objetivo]
Use \texttt{std::endl} apenas quando a mensagem precisa estar no disco
\emph{agora} --- a última linha antes de um aborto previsto, por exemplo. Para
qualquer outra coisa, \texttt{'\textbackslash n'}. O \texttt{Logger} do
subprojeto faz a escolha oposta de propósito, e a
\cref{sec:log-idempotente} explica por quê.
\end{regra}

\paragraph{O conselho que se inverteu.} O \texttt{README.md} do subprojeto
apresenta o \texttt{ostringstream} como ``concatenação eficiente (evita $N$
alocações do \texttt{+} com \texttt{string})''. A medição diz o contrário: o
\texttt{std::string} com \texttt{+=} foi 1,8 vez \emph{mais rápido}. A
explicação é que o conselho envelheceu junto com a implementação. Ele era
verdadeiro contra \texttt{a + b + c}, que constrói temporários; contra
\texttt{+=}, não é --- \texttt{std::string} cresce com capacidade
geométrica, e o \emph{amortizado} de \texttt{+=} é uma cópia de bytes, enquanto
cada \texttt{<<} do \emph{stream} atravessa a maquinaria de formatação,
\emph{locale} e \emph{sentry}. O valor do \texttt{ostringstream} é a
\emph{formatação} (\texttt{setw}, \texttt{setprecision}, tabelas alinhadas),
não a velocidade de emenda.

\paragraph{O conselho que encolheu.} O terceiro é
\texttt{std::ios\_base::sync\_with\_stdio(false)}, que a folclore das maratonas
de programação vende como ganho de ordem de grandeza. Medido aqui, com 300 mil
linhas em \texttt{std::cout} redirecionado para \texttt{/dev/null}:

\begin{lstlisting}[style=bt, language=bash,
    caption={A mesma carga com e sem sincronização com o \texttt{stdio} do C.}]
sync=true  : 12,3 ms
sync=false : 10,1 ms      ← 1,2×, não 10×
\end{lstlisting}

Não é que a folclore esteja errada; é que ela mede outra coisa. O ganho grande
aparece em \emph{leitura} caractere a caractere com \texttt{std::cin}, em que a
sincronização custa uma ida ao \emph{buffer} do C por caractere. Na escrita em
lote, o tempo é dominado pela formatação, e a sincronização some no ruído.

\begin{regra}[um número de \emph{benchmark} pertence a uma carga, não a uma API]
As três medições acima ensinam o mesmo método. Antes de repetir um conselho de
desempenho, pergunte três coisas: \emph{qual carga} foi medida, \emph{em que
versão} da biblioteca, e \emph{o que o compilador pode ter apagado}. Um conselho
que não responde às três é folclore --- inclusive quando é verdadeiro.
\end{regra}

\subsubsection{Redirecionamento: \texttt{rdbuf} e o que ele não captura}

Todo \emph{stream} é uma fachada de formatação sobre um \texttt{std::streambuf},
que é quem de fato move os bytes. Trocar o \emph{buffer} redireciona a saída sem
tocar em nenhuma das chamadas que escrevem nela --- mecanismo que a
\cref{sec:file-io} apresenta e que a \cref{sec:io-streams} detalha, com o
\texttt{RedirectGuard} e o \texttt{TeeBuf} do subprojeto transcritos. O que cabe
acrescentar aqui são duas armadilhas de tempo de vida, ambas da natureza deste
capítulo: erros que compilam.

A primeira explica por que o \texttt{RedirectGuard} declara \emph{deletados} o
construtor de cópia e a atribuição. Copiá-lo faria dois objetos restaurarem o
mesmo \emph{buffer}, e o segundo restauraria sobre um estado já restaurado ---
exatamente a razão pela qual \texttt{unique\_ptr} também não é copiável
(\cref{sec:value-semantics}). Um guarda de recurso é, por construção,
\emph{move-only}.

\begin{armadilha}[o \emph{buffer} redirecionado não pode morrer antes do \emph{stream}]
\texttt{std::cout.rdbuf(arquivo.rdbuf())} guarda em \texttt{std::cout} um
ponteiro para um \texttt{streambuf} que pertence ao \texttt{ofstream}. Se o
\texttt{ofstream} sair de escopo antes da restauração, \texttt{std::cout} fica
com um ponteiro pendurado --- e o próximo \texttt{<<} é \gls{ub}, possivelmente
milhares de linhas depois. O comentário em maiúsculas no
\texttt{demo\_logging.hpp} (``IMPORTANTE: restaurar antes de qualquer
\texttt{fmt::print} ou destruição'') existe por isso. Com o
\texttt{RedirectGuard}, a ordem correta é consequência da ordem de declaração:
o arquivo é declarado antes do guarda, e destruído depois dele.
\end{armadilha}

Há um segundo detalhe, e este é o que mais surpreende num repositório em que
quase toda saída passa por \texttt{fmt::print}. O experimento: redirecionar
\texttt{std::cout} para um arquivo e escrever por três vias diferentes.

\begin{lstlisting}[style=bt, language=C++,
    caption={O experimento --- três escritas, um redirecionamento.}]
std::ofstream f("/tmp/redir_test.txt");
std::streambuf* orig = std::cout.rdbuf(f.rdbuf());
std::cout   << "via std::cout\n";
fmt::print("via fmt::print\n");
std::printf("via printf\n");
std::cout.rdbuf(orig);
\end{lstlisting}

\begin{lstlisting}[style=bt, language=bash,
    caption={Resultado medido: só uma das três foi capturada.}]
--- conteudo do arquivo ---
  via std::cout

--- o que continuou saindo no terminal ---
via fmt::print
via printf
\end{lstlisting}

\begin{armadilha}[trocar o \texttt{rdbuf} de \texttt{cout} não redireciona \texttt{fmt::print}]
\texttt{std::cout.rdbuf()} mexe em um objeto C++; \texttt{fmt::print} e
\texttt{printf} escrevem no \emph{descritor} 1, sem passar por ele. Quem captura
saída para teste trocando o \emph{buffer} de \texttt{std::cout} capturará
exatamente nada num programa que usa \texttt{fmt}. Para redirecionar de verdade
é preciso descer um nível: \texttt{freopen} (nível C, global e difícil de
desfazer) ou \texttt{dup2} (nível POSIX, reversível guardando o \emph{fd}
antigo com \texttt{dup}). É por isso que o \texttt{Logger} do subprojeto recebe
o \emph{stream} de destino como parâmetro em vez de confiar em
redirecionamento: injetar a dependência é mais simples do que interceptá-la.
\end{armadilha}

Na direção oposta está o \texttt{TeeBuf} do subprojeto: em vez de trocar um
destino por outro, escreve nos dois, implementando em C++ o que a \emph{shell}
faria com \texttt{./prog 2>\&1 | tee log.txt}. Ele é um \texttt{std::streambuf}
próprio, e a \cref{sec:io-streams} o transcreve --- vale antecipar apenas que o
contrato mínimo de um \texttt{streambuf} de saída são dois métodos,
\texttt{overflow} e \texttt{xsputn}, e que ele herda a mesma armadilha de tempo
de vida: o \texttt{TeeBuf} aponta para dois \emph{buffers} que precisam
sobreviver a ele.

\subsubsection{O \emph{Logger}, e a lição de idempotência}
\label{sec:log-idempotente}

O \texttt{Logger} do subprojeto reúne o capítulo inteiro num objeto de sessenta
linhas: níveis (\texttt{TRACE} a \texttt{FATAL}), \emph{timestamp} montado com
\texttt{ostringstream} e \texttt{put\_time}, cor ANSI no terminal e texto limpo
no arquivo, \texttt{std::mutex} para ser utilizável de várias \emph{threads}, e
\texttt{ERROR}/\texttt{FATAL} desviados para \texttt{std::cerr}. Esse desvio é
verificável na hora, e vale fazê-lo: rodando com \texttt{1> /dev/null}, restam
na tela apenas as duas últimas linhas; com \texttt{2> /dev/null}, restam todas
as outras.

Há, porém, uma observação que só aparece rodando o programa duas vezes. O
arquivo é aberto com \texttt{std::ios::app}, e o demo imprime o conteúdo
\emph{inteiro} do arquivo ao final:

\begin{lstlisting}[style=bt, language=bash,
    caption={Sétima execução consecutiva de \texttt{dist/bin/file\_io}: a saída
             do demo cresce a cada rodada.}]
  conteúdo:
    [11:33:57.270] [INFO ] Logger iniciado → /tmp/cpp_academy_logger.log
    ...
    [12:44:22.471] [INFO ] Logger iniciado → /tmp/cpp_academy_logger.log
    ...
    [14:24:39.785] [INFO ] Logger iniciado → /tmp/cpp_academy_logger.log
    ...                                     ← oito sessões acumuladas
\end{lstlisting}

Para um \emph{log} de verdade, \texttt{ios::app} é a escolha certa --- perder o
histórico a cada reinício seria pior. Para um \emph{demo}, é um defeito: a
saída do programa depende de quantas vezes ele já rodou, o que impede compará-la
com uma saída de referência e, portanto, testá-la.

\begin{regra}[um programa de demonstração deve ser idempotente]
Rodar duas vezes tem de produzir a mesma saída. Onde o assunto exige estado
persistente, o programa deve começar limpando o que ele mesmo criou. Este
repositório aprendeu a regra do jeito difícil: ao habilitar os alvos que antes
não eram compilados, o demo de \texttt{src/database/} passou na primeira
execução e falhou na segunda com violação de restrição \texttt{UNIQUE}, porque
inseria as mesmas linhas de novo. A correção --- apagar a tabela antes de
popular --- é de uma linha; achar o problema exigiu \emph{rodar duas vezes},
que é justamente o que ninguém faz com um alvo que acabou de compilar.
\end{regra}

\subsubsection{\texttt{std::filesystem}: a API dupla e a corrida}

O \cabecalho{filesystem} do C++17 fecha a seção. A \cref{sec:file-io} já
apresenta o \texttt{fs::path} e os iteradores de diretório; o que cabe aqui,
enquanto o assunto ainda é tratamento de erro, são duas propriedades da API.

A primeira é que \textbf{toda função tem duas assinaturas}: uma que lança
\texttt{fs::filesystem\_error} e outra que recebe um \texttt{std::error\_code\&}
e não lança. O demo executa as duas sobre o mesmo arquivo inexistente, e a
diferença de qualidade da mensagem é visível:

\begin{lstlisting}[style=bt, language=bash,
    caption={\texttt{dist/bin/file\_io}, demo 5, seção 5.}]
  exceção: filesystem error: cannot get file size:
           No such file or directory [/tmp/nao_existe_absolutamente.xyz]
  path1:   /tmp/nao_existe_absolutamente.xyz
  code:    No such file or directory
  error_code: No such file or directory (valor=2)
\end{lstlisting}

A exceção carrega o \emph{caminho} --- \texttt{path1()} e, em operações de dois
argumentos, \texttt{path2()} --- que o \texttt{error\_code} sozinho não tem. O
critério de escolha é a frequência esperada do erro: um arquivo de configuração
que deveria existir e não existe é excepcional; percorrer um diretório em que
arquivos podem sumir no meio do caminho, não.

A segunda propriedade é uma armadilha de concorrência que a API não tem como
evitar.

\begin{armadilha}[\texttt{exists()} responde sobre o passado]
\texttt{if (fs::exists(p)) abrir(p);} é uma corrida: entre a pergunta e a
abertura, outro processo pode apagar, criar ou substituir \texttt{p} por um
\emph{link} simbólico. O padrão tem nome --- TOCTOU, \emph{time of check to time
of use} --- e é uma classe inteira de vulnerabilidades de segurança. O remédio é
não perguntar: tente a operação e trate a falha, que é o que o par
\texttt{ifstream} + \texttt{is\_open()} já faz naturalmente. Use
\texttt{exists()} para relatar, nunca para decidir.
\end{armadilha}

Vale registrar ainda que \texttt{fs::permissions} manipula os bits POSIX com
\texttt{perm\_options::add} e \texttt{::remove}, e que o demo os usa para
remover a permissão de escrita do dono --- e depois restaurá-la, sem o que o
\texttt{fs::remove} final falharia. O detalhe é pequeno e ensina algo geral: um
demo que altera permissões precisa desfazer o que fez, pela mesma razão da
regra de idempotência acima.

\begin{aviso}
O \texttt{README.md} do \projeto{file\_io} ainda instrui a ``ativar o
subprojeto em \texttt{src/meson.build}'' descomentando um \texttt{subdir}. Essa
instrução está obsoleta: \textbf{não há mais nenhum \texttt{subdir()}
comentado} --- os 57 subprojetos são construídos, e a compilação instala 98
binários em \texttt{dist/bin}, sem um \emph{warning} sob
\texttt{-Wall -Wextra -Wpedantic}. A mudança não foi cosmética: ao habilitar os
alvos que estavam desligados, apareceram quatro \texttt{\#include} faltando ---
entre eles o \texttt{<fmt/core.h>} de \texttt{src/aws/s3/main.cpp} --- e dois
bugs de ABI em \texttt{src/asm\_course/} que estavam ali havia meses.
\textbf{Um alvo que não é construído é um alvo que apodrece}: o compilador é o
único leitor que nunca se distrai, e desligá-lo de um diretório equivale a
parar de revisar aquele diretório.
\end{aviso}

% =============================================================================
\subsection{O modelo de objeto}
\label{sec:oop}
% =============================================================================

Herança e \texttt{virtual} aparecem em quase todo subprojeto do repositório ---
\projeto{design-of-pattern}, \projeto{plugin\_registry}, \projeto{ecs} --- mas
sempre \emph{usados}, nunca \emph{explicados}. O subprojeto \projeto{oop} fecha
essa lacuna: o que o compilador gera por baixo do \texttt{virtual}, e os quatro
modos clássicos de o polimorfismo desaparecer sem que ninguém perceba.

\subsubsection{Herança}

Das três formas de herdar, apenas a pública cria substituibilidade
(\cref{tab:heranca}). As outras duas são composição com sintaxe diferente --- e
composição costuma ser a escolha melhor, com a vantagem de não convidar ninguém
a fazer \emph{upcast}.

\begin{table}[htbp]
\centering
\caption{As três formas de herança.}
\label{tab:heranca}
\footnotesize
\begin{tabular}{@{}lll@{}}
\toprule
\textbf{Forma} & \textbf{Significado} & \textbf{Converte para a base?} \\
\midrule
\texttt{public}    & ``é um''                        & sim \\
\texttt{protected} & ``é implementado em termos de'' & não \\
\texttt{private}   & ``é implementado em termos de'' & não \\
\bottomrule
\end{tabular}
\end{table}

A ordem de construção é base, depois membros na ordem de \emph{declaração},
depois o corpo do construtor; a destruição é o inverso exato.

\begin{armadilha}[declarar um membro esconde \emph{todos} os \emph{overloads} do mesmo nome]
\begin{lstlisting}[style=bt, language=C++, frame=none,
                   xleftmargin=0pt, xrightmargin=0pt]
struct Base    { void f(int); void f(const std::string&); };
struct Ruim : Base { void f(double); };          // escondeu os dois

Ruim r;
r.f(42);          // chama f(double) — 42 vira 42.0, em silêncio
r.f("texto");     // nem compila
\end{lstlisting}
A busca de nome para assim que encontra o nome, \emph{antes} da resolução de
sobrecarga: os \emph{overloads} da base nem chegam a ser candidatos. A correção
é uma linha --- \texttt{using Base::f;}.
\end{armadilha}

\subsubsection{Funções virtuais e a \emph{vtable}}

O mecanismo é conhecido, mas convém dizê-lo com precisão: cada classe
polimórfica ganha uma \gls{vtable} --- um vetor de ponteiros de função --- e cada
\emph{objeto} ganha um \emph{vptr} apontando para a tabela da sua classe real.
A chamada \texttt{p->f()} vira, aproximadamente,
\texttt{(*p->vptr[índice\_de\_f])(p)}. O subprojeto mede as consequências:

\begin{lstlisting}[style=bt, language=bash,
    caption={Saída de \texttt{src/oop/}: a \emph{vtable} é por classe, o
             \emph{vptr} é por objeto.}]
sizeof(SemVirtual) = 8 bytes (2 ints)
sizeof(ComVirtual) = 16 bytes (2 ints + vptr + padding)
vptr de dois Círculos são iguais? true    ← vtable é POR CLASSE
vptr de Círculo == vptr de Forma?  false  ← vptr é POR OBJETO
\end{lstlisting}

O custo em memória é um ponteiro por objeto. O custo maior, quase sempre, é
outro: a chamada virtual \textbf{não pode sofrer \emph{inline}}, e com isso se
perdem as otimizações que dependiam de ver o corpo da função.

\paragraph{\texttt{override} e o bug silencioso.} Uma assinatura ligeiramente
diferente não sobrescreve nada --- cria um \emph{overload} novo, e as chamadas
pela base continuam indo para a base:

\begin{lstlisting}[style=bt, language=C++,
    caption={O bug que \texttt{override} previne, em
             \texttt{src/oop/demo\_virtual.hpp}.}]
struct Sensor {
    virtual double ler(int /*canal*/) const { fmt::print("      Sensor::ler(int)\n"); return 0.0; }
    virtual ~Sensor() = default;
};

struct SensorQuebrado : Sensor {
    // Repare no `double` em vez de `int`: assinatura diferente.
    // Isto NÃO sobrescreve nada — cria um novo overload.
    double ler(double canal) const { fmt::print("      SensorQuebrado::ler(double)\n");
                                     return canal; }
};

struct SensorCerto : Sensor {
    // `override` faz o compilador VERIFICAR que existe um virtual com esta
    // assinatura na base. Se não existir, é erro de compilação.
    double ler(int canal) const override { fmt::print("      SensorCerto::ler(int)\n");
                                           return canal; }
};
\end{lstlisting}

\begin{regra}[\texttt{override} sempre; \texttt{final} quando a classe não é para herdar]
\texttt{override} custa zero em tempo de execução e transforma um bug silencioso
em erro de compilação. \texttt{final} vai além: impede herdar a classe ou
sobrescrever o método, e permite ao compilador \emph{devirtualizar} chamadas ---
o que devolve a possibilidade de \emph{inline}.
\end{regra}

\paragraph{Destrutor virtual.} \texttt{delete ptr\_base} com destrutor não
virtual é \gls{ub}: o destrutor do derivado simplesmente não roda, e o que ele
liberaria vaza. O GCC recusa com \texttt{-Wdelete-non-virtual-dtor} e o
AddressSanitizer aborta o programa --- razão pela qual a linha correspondente
está comentada no subprojeto, de propósito.

\paragraph{Virtual dentro do construtor.} Durante o construtor da base, o
\emph{vptr} ainda é o da base: a chamada não desce para o derivado. Faz sentido
--- os membros do derivado ainda não existem --- e é exatamente por isso que a
construção em duas fases (``construtor chama \texttt{inicializar()} virtual'')
não funciona em C++.

\begin{armadilha}[argumento padrão em função virtual]
\begin{lstlisting}[style=bt, language=C++, frame=none,
                   xleftmargin=0pt, xrightmargin=0pt]
struct BaseArg {
    virtual void f(int x = 10) const { fmt::print("      BaseArg::f(x={})\n", x); }
    virtual ~BaseArg() = default;
};
struct DerivadaArg : BaseArg {
    void f(int x = 20) const override { fmt::print("      DerivadaArg::f(x={})\n", x); }
};
\end{lstlisting}
O \emph{corpo} é escolhido dinamicamente, mas o \emph{argumento padrão} é
escolhido estaticamente, pelo tipo do ponteiro. Chamado por
\texttt{Derivada d; d.f()} imprime \texttt{x=20}; chamado por
\texttt{Base\& b = d; b.f()} imprime \texttt{x=10} --- corpo do derivado, valor
padrão da base. A solução é o idioma NVI: um método público não virtual aplica
o padrão e delega a um virtual privado.
\end{armadilha}

\subsubsection{\emph{Object slicing}}

Copiar um derivado para uma variável do tipo base não é conversão: é
\textbf{amputação}. O que não cabe em \texttt{sizeof(Base)} é descartado, e com
ele o \emph{vptr} --- ou seja, o polimorfismo.

\begin{lstlisting}[style=bt, language=C++,
    caption={Os tipos de \texttt{src/oop/demo\_slicing.hpp}: o estado extra é
             o que some no corte.}]
struct Animal {
    std::string nome;

    explicit Animal(std::string n) : nome(std::move(n)) {}
    virtual ~Animal() = default;

    virtual std::string falar() const { return "..."; }
    virtual std::string descrever() const { return fmt::format("Animal '{}'", nome); }
};

struct Cachorro : Animal {
    int  numero_de_truques;         // ← estado extra, some no corte

    Cachorro(std::string n, int truques)
        : Animal(std::move(n)), numero_de_truques(truques) {}

    std::string falar() const override { return "Au!"; }
    std::string descrever() const override {
        return fmt::format("Cachorro '{}' com {} truques", nome, numero_de_truques);
    }
};
\end{lstlisting}

O caso que aparece em produção não é a atribuição explícita: é o
\textbf{parâmetro}. A diferença entre \texttt{void f(Animal a)} e \texttt{void
f(const Animal\& a)} é um caractere, e nenhum \emph{warning} separa as duas.

\begin{armadilha}[\texttt{std::vector<Base>} é sempre errado]
O container reserva \texttt{sizeof(Base)} por posição. Não existe espaço físico
para um derivado maior --- não é uma questão de política, é de aritmética:
\begin{lstlisting}[style=bt, language=bash, frame=none,
                   xleftmargin=0pt, xrightmargin=0pt]
sizeof(Animal)   = 40 bytes
sizeof(Cachorro) = 48 bytes
→ dois Cachorros entraram no vector<Animal>, dois Animals saíram
\end{lstlisting}
A correção é guardar indireção: \texttt{std::vector<std::unique\_ptr<Animal>>}.
\end{armadilha}

O corte pode ser tornado incompilável, e vale escolher a receita conforme a
intenção: base abstrata (um \texttt{= 0}) é a mais forte, porque nem instancia;
cópia \texttt{= delete} na base tira a cópia também dos derivados; cópia
\texttt{protected} deixa derivados copiarem entre si e impede o mundo externo de
cortar.

\subsubsection{Herança múltipla e o diamante}

Herdar de várias \emph{interfaces} puras --- só virtuais puros, sem estado ---
não tem o problema do diamante, porque não há dados a duplicar. Mas há um efeito
de \emph{layout} que surpreende: os subobjetos ficam em \emph{offsets}
diferentes, e converter o ponteiro \textbf{ajusta o endereço}.

\begin{lstlisting}[style=bt, language=bash,
    caption={Saída de \texttt{src/oop/}: o mesmo objeto, três endereços.}]
&b            = 0x7ffc2c9a4700
(Desenhavel*) = 0x7ffc2c9a4700
(Salvavel*)   = 0x7ffc2c9a4708   ← mesmo objeto, endereço diferente
\end{lstlisting}

Comparar ponteiros de bases diferentes, portanto, não responde se são o mesmo
objeto. Quem responde é \texttt{dynamic\_cast<void*>}, que devolve o endereço do
objeto mais derivado.

Quando as bases têm estado, o diamante aparece: sem \texttt{virtual}, a classe
final contém \emph{duas} cópias da base comum, com estados independentes, e a
conversão para ela é ambígua --- nem compila. Com herança virtual, contém uma,
e o \texttt{sizeof} cai (96 para 64 bytes, no exemplo do subprojeto).

\begin{armadilha}[quem inicializa a base virtual é a classe \emph{mais derivada}]
Numa hierarquia em diamante com herança virtual, as chamadas ao construtor da
base virtual escritas nas classes intermediárias são simplesmente
\textbf{ignoradas}. Quem escolhe com que argumentos a base virtual é construída
é a classe mais derivada --- e isso significa que toda classe concreta da
hierarquia precisa saber inicializar uma base da qual talvez nem se considere
responsável.
\end{armadilha}

Quando a variação é conhecida em compilação, nada disso é necessário: o
\gls{crtp} dá polimorfismo estático sem \emph{vtable} e sem \emph{vptr}, e
\emph{mixins} vazios desaparecem por otimização de base vazia
(\cref{sec:templates}).

\subsubsection{RTTI e o seu custo}
\label{sec:oop-rtti}

\texttt{typeid} aplicado a um objeto \textbf{polimórfico} é avaliado em tempo de
execução e devolve o tipo \emph{dinâmico}; aplicado a um não polimórfico, é
resolvido em compilação e devolve o tipo \emph{estático}. A diferença entre as
duas situações é uma função virtual na base.

\begin{lstlisting}[style=bt, language=bash,
    caption={Saída de \texttt{src/oop/}: o nome cru vem no formato do
             Itanium ABI~\cite{itaniumabi} e precisa de
             \texttt{abi::\_\_cxa\_demangle}.}]
typeid(c).name() cru      = N9demo_rtti7CirculoE
demanglado                = demo_rtti::Circulo
typeid(f) via Forma&      = demo_rtti::Circulo          ← dinâmico
não-polimórfica via base& = demo_rtti::NaoPolimorfica   ← estático
\end{lstlisting}

\texttt{typeid} testa igualdade \emph{exata} de tipo; \texttt{dynamic\_cast}
aceita derivados. São perguntas diferentes: um quadrado que herda de retângulo
falha no primeiro teste e passa no segundo.

A medição do subprojeto, sobre 200\,000 iterações de cada operação, ordena os
custos: chamada virtual é a mais barata, comparação de \texttt{typeid} vem
depois, e \texttt{dynamic\_cast} é a mais cara por boa margem --- porque
percorre a árvore de herança e não é $O(1)$.

\begin{regra}[uma cadeia de \texttt{if (dynamic\_cast<X>)} é uma função virtual que não foi escrita]
As três alternativas, em ordem de preferência: (a) um método virtual, e o objeto
responde por si; (b) o padrão \emph{Visitor}, que faz despacho duplo sem RTTI
(\cref{sec:patterns}); (c) \texttt{std::variant} com \texttt{std::visit}, quando
a hierarquia é fechada --- sem herança nenhuma, e com verificação de
exaustividade pelo compilador (\cref{sec:vocabulary}).
\end{regra}

% =============================================================================
\subsection{Tratamento de erro}
\label{sec:errors}
% =============================================================================

O subprojeto \projeto{errors} trata do mecanismo das exceções, dos quatro níveis
de garantia que toda função oferece --- escolhendo-os ou não --- e das duas
alternativas baseadas em valor. Ele compila com C++23 por causa de
\texttt{std::expected} (\cref{sec:ambiente}).

A decisão entre exceção e valor de retorno não é questão de gosto. O subprojeto
mede:

\begin{lstlisting}[style=bt, language=bash,
    caption={Saída de \texttt{src/errors/}: 20\,000 iterações.}]
sem exceção        :      12 µs
lançando+capturando:   12716 µs      ← ~1000×
\end{lstlisting}

Exceções são \textbf{zero-cost no caminho feliz} e caríssimas no caminho de
erro. É isso, e não a estética, que decide quando usá-las: para o que é
excepcional e raro, exceção; para o que é esperado e frequente, valor.

\subsubsection{O mecanismo, e por que ele é o mesmo assunto que RAII}

Um \texttt{throw} faz quatro coisas: copia ou move o objeto lançado para uma
área controlada pelo \emph{runtime}; desenrola a pilha até o primeiro
\emph{handler} compatível; \textbf{roda todo destrutor de objeto local} durante o
desenrolamento; executa o \emph{handler}. O terceiro passo é o que torna
\gls{raii} e exceções indissociáveis:

\begin{lstlisting}[style=bt, language=bash,
    caption={Saída de \texttt{src/errors/}: três destrutores rodaram sem uma
             linha de limpeza escrita à mão.}]
+ nível-1 construído
+ nível-2 construído
+ nível-3 construído
nível-3 vai lançar...
- nível-3 DESTRUÍDO (durante o unwinding)
- nível-2 DESTRUÍDO (durante o unwinding)
- nível-1 DESTRUÍDO (durante o unwinding)
handler alcançado: erro originado no nível 3
\end{lstlisting}

Os \emph{handlers} são testados \textbf{em ordem}, de cima para baixo, e o
primeiro compatível vence --- não há ``melhor correspondência'' como na
resolução de sobrecarga. Por isso o \emph{handler} mais derivado vem primeiro.

\begin{armadilha}[capturar por valor \emph{fatia} a exceção]
\texttt{catch (std::runtime\_error e)} sofre \gls{slicing} exatamente como
qualquer outra cópia para o tipo base (\cref{sec:oop}): o objeto derivado chega
degradado, e o \texttt{what()} vira genérico. O GCC pega com
\texttt{-Wcatch-value}. Lançar, por outro lado, é sempre por valor:
\texttt{throw T(...)}, nunca \texttt{throw new T(...)} --- este último vaza.
\end{armadilha}

A mesma armadilha reaparece no repasse. \texttt{throw;} relança \emph{o mesmo}
objeto, com tipo dinâmico intacto; \texttt{throw e;} lança uma \textbf{cópia}
cortada para o tipo estático da variável do \texttt{catch}:

\begin{lstlisting}[style=bt, language=bash,
    caption={Saída de \texttt{src/errors/}: a diferença de um caractere.}]
com `throw;'     → chegou como ErroDeRede: [HTTP 404] não encontrado
com `throw e;'   → chegou degradado como std::exception: 'std::exception'
\end{lstlisting}

\subsubsection{As quatro garantias}

Toda função oferece uma das garantias da \cref{tab:garantias}, tenha o autor
pensado nisso ou não. A escolha razoável é mirar na básica sempre, e na forte
quando o custo permitir.

\begin{table}[htbp]
\centering
\caption{As quatro garantias de segurança em face de exceções.}
\label{tab:garantias}
\footnotesize
\begin{tabular}{@{}ll@{}}
\toprule
\textbf{Garantia} & \textbf{Promessa} \\
\midrule
nenhuma  & nada --- pode vazar ou corromper \\
básica   & não vaza, invariantes valem, valores não especificados \\
forte    & ou completou, ou nada mudou (\emph{commit-or-rollback}) \\
\emph{nothrow} & não lança, ponto \\
\bottomrule
\end{tabular}
\end{table}

Subir de nível costuma custar menos do que parece --- às vezes apenas
\textbf{reordenar as operações}, deixando o que pode lançar antes do que muda
estado:

\begin{lstlisting}[style=bt, language=C++,
    caption={A mesma função, com garantia diferente.}]
void ruim(const std::string& k) {
    ++contador_;                                    // estado já mudou
    itens_.push_back(std::make_unique<Item>(k));    // e isto pode lançar
}

void bom(const std::string& k) {
    auto novo = std::make_unique<Item>(k);          // o arriscado primeiro
    itens_.push_back(std::move(novo));
    ++contador_;                                    // o nothrow por último
}
\end{lstlisting}

Para a garantia forte, o instrumento é o mesmo \emph{copy-and-swap} da
\cref{sec:value-semantics}: faça o trabalho arriscado numa cópia e, só quando
tudo tiver dado certo, troque --- porque \texttt{swap} é \emph{nothrow}.

\begin{armadilha}[destrutor que lança durante o \emph{unwinding} chama \classe{std::terminate}]
Se um destrutor lançar enquanto outra exceção está em voo, há duas exceções
simultâneas e o padrão responde encerrando o programa. Por isso todo destrutor é
implicitamente \texttt{noexcept}. O padrão correto é engolir e registrar o erro
no destrutor, e oferecer um \texttt{fechar()} público para quem quiser tratá-lo
de verdade --- é assim que \texttt{std::ofstream} se comporta.
\end{armadilha}

\begin{regra}[\texttt{noexcept} é promessa, não impedimento]
\texttt{noexcept} não impede a função de lançar: promete que ela não lança.
Quebrar a promessa não produz \emph{unwinding} nem \emph{handler} ---
\texttt{std::terminate} direto. Em troca, a biblioteca otimiza: é o que faz
\texttt{std::vector} mover em vez de copiar ao realocar. Marque
\texttt{noexcept} em construtores e atribuições de \emph{move}, em
\texttt{swap}, em destrutores e em \emph{getters} triviais.
\end{regra}

\subsubsection{A hierarquia padrão}

A pergunta que separa os dois grandes ramos é uma só: \emph{isto poderia ter
sido evitado por quem chamou, verificando antes?}

\begin{itemize}
    \item \textbf{Sim} $\rightarrow$ \classe{std::logic\_error} e derivados
          (\texttt{invalid\_argument}, \texttt{domain\_error},
          \texttt{length\_error}, \texttt{out\_of\_range}). É bug. Não capture
          em produção; corrija o código.
    \item \textbf{Não} $\rightarrow$ \classe{std::runtime\_error} e derivados
          (\texttt{range\_error}, \texttt{overflow\_error},
          \texttt{system\_error}). É o mundo. Trate, tente de novo, faça
          \emph{fallback}.
\end{itemize}

Uma exceção customizada deve derivar de \texttt{std::runtime\_error}, que já
gerencia a \emph{string} do \texttt{what()} --- e essa \emph{string} precisa
sobreviver ao objeto lançado, o que é a razão de não se poder montá-la na hora
dentro do \texttt{what()} e devolver \texttt{.c\_str()} de um temporário.

Vale também dar ao seu domínio uma \textbf{raiz própria}:

\begin{lstlisting}[style=bt, language=C++,
    caption={Uma raiz por domínio permite capturar tudo o que é seu sem pegar o
             que não é.}]
struct ErroDaAcademia : std::runtime_error { using std::runtime_error::runtime_error; };
struct ErroDeParse    : ErroDaAcademia     { using ErroDaAcademia::ErroDaAcademia; };
struct ErroDeRede     : ErroDaAcademia     { int codigo; /* ... */ };
\end{lstlisting}

Para acrescentar contexto sem apagar a causa, \texttt{std::throw\_with\_nested}
na subida e \texttt{std::rethrow\_if\_nested} na leitura produzem o rastro
encadeado:

\begin{lstlisting}[style=bt, language=bash,
    caption={Saída de \texttt{src/errors/}: cada camada acrescenta contexto.}]
falha ao iniciar o servidor na porta 443
  não foi possível carregar o certificado TLS
    arquivo '/etc/ssl/cert.pem' não encontrado
\end{lstlisting}

\begin{regra}[\texttt{catch(...)} é legítimo em exatamente três lugares]
No \texttt{main}, na função de entrada de uma \emph{thread} --- onde uma
exceção que escapa chama \texttt{std::terminate}, e o remédio é
\texttt{std::exception\_ptr} --- e numa fronteira com código C, que não sabe
desenrolar pilha de C++. Em qualquer outro ponto, \texttt{catch(...)} está
escondendo um bug.
\end{regra}

\subsubsection{Erros como valores}

\classe{std::error\_code} é um par (valor inteiro, categoria) que identifica um
erro de forma extensível, sem alocar e sem lançar. A distinção que o torna
interoperável é entre \texttt{error\_code} --- o que \emph{aconteceu},
específico da plataforma ou da biblioteca --- e \texttt{error\_condition} --- o
que aquilo \emph{significa}, portável. Uma categoria própria que implemente
\texttt{default\_error\_condition} passa a ser comparável com
\texttt{std::errc}, e é assim que Asio, \texttt{<filesystem>} e o código do
usuário coexistem sem conhecer uns aos outros.

\begin{armadilha}[a polaridade do \texttt{bool} é invertida entre as duas formas]
Com \texttt{std::error\_code}, \texttt{if (ec)} significa \textbf{houve erro}.
Com \texttt{std::expected}, \texttt{if (r)} significa \textbf{sucesso}. As duas
convenções coexistem no mesmo programa, e a confusão entre elas produz código
que inverte o tratamento de erro sem que nada deixe de compilar.
\end{armadilha}

\texttt{std::expected<T, E>} (C++23) resolve o que faltava às outras formas:
\texttt{optional<T>} diz que não há valor mas não diz por quê;
\texttt{error\_code} diz o porquê mas não carrega o valor.

\begin{lstlisting}[style=bt, language=C++,
    caption={\texttt{std::expected} em
             \texttt{src/errors/demo\_expected.hpp}.}]
std::expected<int, Falha> para_inteiro(std::string_view s) noexcept {
    if (s.empty()) return std::unexpected{Falha::vazio};

    int valor = 0;
    const auto [ptr, ec] = std::from_chars(s.data(), s.data() + s.size(), valor);

    if (ec == std::errc::result_out_of_range) return std::unexpected{Falha::fora_do_intervalo};
    if (ec != std::errc{} || ptr != s.data() + s.size())
        return std::unexpected{Falha::caractere_invalido};

    return valor;                    // sucesso: conversão implícita para expected
}
\end{lstlisting}

O ganho real, porém, é o encadeamento monádico: a primeira falha curto-circuita
todas as etapas seguintes, e nenhum \texttt{if} aparece no meio do
\emph{pipeline}.

\begin{lstlisting}[style=bt, language=C++,
    caption={O \emph{pipeline} de \texttt{src/errors/demo\_expected.hpp}.}]
    auto pipeline = [](std::string_view entrada) {
        return para_inteiro(entrada)              // expected<int, Falha>
            .and_then(exigir_positivo)            // só roda se deu certo
            .transform([](int v) { return v * 2; })   // mapeia só o sucesso
            .transform([](int v) { return fmt::format("resultado={}", v); });
    };
\end{lstlisting}

\begin{table}[htbp]
\centering
\caption{As quatro operações monádicas de \texttt{std::expected}.}
\label{tab:monadicas}
\footnotesize
\begin{tabular}{@{}llll@{}}
\toprule
\textbf{Operação} & \textbf{Recebe} & \textbf{Devolve} & \textbf{Papel} \\
\midrule
\texttt{and\_then(f)}        & \texttt{T} & \texttt{expected<U,E>} & encadeia \\
\texttt{transform(f)}        & \texttt{T} & \texttt{U}             & mapeia o sucesso \\
\texttt{or\_else(f)}         & \texttt{E} & \texttt{expected<T,G>} & recupera \\
\texttt{transform\_error(f)} & \texttt{E} & \texttt{G}             & traduz o erro \\
\bottomrule
\end{tabular}
\end{table}

\begin{armadilha}[\texttt{std::expected} não é \texttt{[[nodiscard]]} no padrão]
Ignorar o retorno de uma função que devolve \texttt{expected} compila sem um
pio --- e o erro desaparece por completo, o que é pior do que a exceção que se
queria evitar. Marque você as suas funções:
\texttt{[[nodiscard]] std::expected<int,Falha> para\_inteiro(...);}
\end{armadilha}

\subsubsection{Escolhendo a forma}

\begin{table}[htbp]
\centering
\caption{As cinco formas de reportar erro, comparadas.}
\label{tab:formas-erro}
\footnotesize
\begin{tabular}{@{}llll@{}}
\toprule
\textbf{Forma} & \textbf{Motivo do erro} & \textbf{Custo} & \textbf{Ignorável?} \\
\midrule
\texttt{optional<T>}          & não             & zero  & sim \\
exceção                       & sim, rico       & alto  & não \\
\emph{out-param} + \texttt{bool} & via bool/enum & zero  & sim, e em silêncio \\
\texttt{error\_code}          & sim, extensível & zero  & sim \\
\texttt{expected<T,E>}        & sim, \textbf{tipado} & zero & sim, se não marcar \\
\bottomrule
\end{tabular}
\end{table}

E os dois critérios que decidem, na prática:

\begin{itemize}
    \item \textbf{Frequência.} Esperado e frequente $\rightarrow$
          \texttt{expected} ou \texttt{error\_code}. Excepcional e raro
          $\rightarrow$ exceção. Bug do chamador $\rightarrow$ \texttt{assert}
          ou \texttt{logic\_error}.
    \item \textbf{Distância.} Se o erro precisa atravessar muitas camadas que
          nada têm a ver com ele, exceção --- propagar \texttt{expected} muda a
          assinatura de todas elas. E em construtor ou operador, exceção é a
          única opção: não há valor de retorno disponível.
\end{itemize}

% =============================================================================
\subsection{Síntese: os erros que o compilador não denuncia}
% =============================================================================

O capítulo inteiro se resume na \cref{tab:sintese-fundamentos}, que reúne as
armadilhas das oito seções e --- a coluna que importa --- o que fazer a
respeito.

\begin{table}[htbp]
\centering
\caption{Os erros silenciosos do capítulo, e a correção de cada um.}
\label{tab:sintese-fundamentos}
\footnotesize
\begin{tabular}{@{}p{0.30\textwidth}p{0.22\textwidth}p{0.36\textwidth}@{}}
\toprule
\textbf{Erro} & \textbf{Seção} & \textbf{Correção} \\
\midrule
\texttt{int t = 3.99;} trunca              & \ref{sec:syntax}
    & inicialize com \texttt{\{\}} \\
lista de inicialização fora de ordem       & \ref{sec:syntax}
    & \texttt{-Wreorder}; siga a declaração \\
\emph{cast} à moda C esconde qual conversão & \ref{sec:casts}
    & \texttt{-Wold-style-cast} \\
escrever em objeto \texttt{const}          & \ref{sec:casts}
    & não use \texttt{const\_cast} para isso \\
ler \texttt{float} por \texttt{int*}       & \ref{sec:casts}
    & \texttt{std::bit\_cast} ou \texttt{memcpy} \\
destrutor vazio suprime os \emph{moves}    & \ref{sec:value-semantics}
    & Rule of 0; ou \texttt{= default} nos cinco \\
\emph{move ctor} sem \texttt{noexcept}     & \ref{sec:value-semantics}
    & marque \texttt{noexcept} \\
\texttt{return std::move(local)}           & \ref{sec:value-semantics}
    & \texttt{return local;} \\
\texttt{move} onde cabia \texttt{forward}  & \ref{sec:value-semantics}
    & \texttt{forward<T>} em tipo deduzido \\
sobrecarregar \texttt{\&\&} ou \texttt{||} & \ref{sec:operators}
    & não sobrecarregue \\
construtor de um argumento sem \texttt{explicit} & \ref{sec:operators}
    & \texttt{explicit} por padrão \\
\texttt{a[i, j]} recompilado em C++23      & \ref{sec:operators}
    & \texttt{-Wcomma-subscript} \\
ciclo de \texttt{shared\_ptr}              & \ref{sec:smart-pointers}
    & \texttt{weak\_ptr} num dos lados \\
\texttt{shared\_ptr<T>(this)} num método   & \ref{sec:smart-pointers}
    & \texttt{enable\_shared\_from\_this} \\
\emph{weak} vivo retém o objeto inteiro    & \ref{sec:smart-pointers}
    & \texttt{shared\_ptr(new)} se for grande \\
\texttt{release()} usado como empréstimo   & \ref{sec:smart-pointers}
    & \texttt{get()}, ou passe \texttt{*p} \\
\texttt{while (!ifs.eof())}                & \ref{sec:streams-raii}
    & ponha a leitura na condição \\
\texttt{gcount()} lido fora do laço        & \ref{sec:streams-raii}
    & consulte logo após cada \texttt{read} \\
\texttt{'─'} entre aspas simples           & \ref{sec:streams-raii}
    & literal de \texttt{string}; UTF-8 tem 3 bytes \\
\texttt{40 - s.size()} dá a volta          & \ref{sec:streams-raii}
    & \texttt{-Wsign-conversion}; converta antes \\
\texttt{cout.rdbuf()} não pega \texttt{fmt::print} & \ref{sec:streams-raii}
    & \texttt{dup2}, ou injete o \emph{stream} \\
\texttt{std::endl} em laço                 & \ref{sec:streams-raii}
    & \texttt{'\textbackslash n'}; 32$\times$ medido \\
\texttt{fs::exists()} antes de abrir       & \ref{sec:streams-raii}
    & tente e trate a falha (TOCTOU) \\
método derivado esconde \emph{overloads}   & \ref{sec:oop}
    & \texttt{using Base::f;} \\
assinatura errada ao ``sobrescrever''      & \ref{sec:oop}
    & \texttt{override} sempre \\
\texttt{delete} por base sem dtor virtual  & \ref{sec:oop}
    & destrutor virtual, ou \texttt{final} \\
\emph{slicing} em parâmetro por valor      & \ref{sec:oop}
    & \texttt{const T\&}; base abstrata \\
argumento padrão em virtual                & \ref{sec:oop}
    & idioma NVI \\
capturar exceção por valor                 & \ref{sec:errors}
    & \texttt{catch (const T\&)}; \texttt{-Wcatch-value} \\
\texttt{throw e;} para repassar            & \ref{sec:errors}
    & \texttt{throw;} \\
retorno \texttt{expected} ignorado         & \ref{sec:errors}
    & \texttt{[[nodiscard]]} \\
\bottomrule
\end{tabular}
\end{table}

Quatro delas --- \emph{slicing}, destrutor não virtual, \emph{move} suprimido e
uso após \emph{move} --- envolvem tempo de vida ou \emph{layout} de objeto, e
para essas o instrumento certo não é a leitura atenta, é o \emph{sanitizer}
(\cref{sec:sanitizers}). Os subprojetos \projeto{value\_semantics},
\projeto{oop} e \projeto{errors} recomendam explicitamente uma passada sob
\texttt{-Db\_sanitize=address,undefined} justamente por isso.
